% Auto-generated by scripts/06_emit_structure.py
% Source: scans 108–150
% Section id: chapter/4
% Phase 3 deliverable. Footnotes (Phase 4), citations (Phase 5),
% and Wade-Giles → Pinyin (Phase 6) are not yet applied here.

\chapter{Dating the Inscriptions: Relative Chronology}
\label{ch:4}

% source: scan 108, printed 88
% intro-echo-dropped: Dating the Inscriptions:
Relative Chronology\footnote[1]{In this book, I attempt to specify the period of every inscription cited when it is important to do so. In some of these cases I am simply quoting the opinion of the editor of the collection in which the piece appears or the opinion of other scholars; but for those touchstone inscriptions for which I believe the period given is firmly established by basic criteria such as ancestral titles, diviner's name (which may not appear in the inscription itself but elsewhere on the same bone or shell), or writing style (when, as in period V, it is unmistakable), the number (i.e., chüan, page, and piece number, or rubbing number; see \ref{ch:3} (see ch. 3), n. 5; is given in italics. The distinction between, say, \pinyinterm{jiabian} 562 (period IIa) and \pinyinterm{jiabian} 2869 (period IIb), therefore, indicates that a greater degree of uncertainty about the date exists in my mind, at least) in the first case. Note that this italic distinction is only employed when the period is specifically mentioned.}

\section[Introduction]{\origsecnum{4.1}Introduction}
\label{sec:chapters-ch04:4.1}

It is the periodization of the inscriptions—how they may be dated with regard
to one another, and with regard to the royal reigns—that permits us to trace the evolution
of \pinyinterm{shang} divination procedures and \pinyinterm{shang} institutions and beliefs. It allows us to view the
late \pinyinterm{shang} historically, as an evolving dynasty and culture. The issues involved in relative
dating are technical and, at times, frustrating, but their understanding is fundamental to
the historian's task.¹
Given the criteria we possess, it is not yet possible to establish a relative date for every
inscribed bone and shell fragment. In some cases, none of our criteria apply; in other cases,
the ones that do are not specific enough to be useful. As a result, many inscriptions have
been dated on subjective grounds. It is not my intention to resolve the disputes that divide
the field; here I wish only to delineate the criteria involved, which, in virtually every case,
stand in need of refinement.\footnote[2]{The major dispute concerns the dating of the RFG inscriptions and hence the existence of alternating Old and New Schools of ritual practice. See \ref{ch:2} (see ch. 2), n. 18. 0 U ] J}
Rules of Thumb
Some rules of thumb will permit the casual student or the researcher pressed for
time to discover the date which other scholars may have assigned a particular inscription.
% source: scan 109, printed 89
First, it should be noted that the inscriptions in Shima Inkyo bokuji sōrui (sec. 3.3.2) are
arranged chronologically within each entry; once an inscription has been located in Sōrui,
it is possible to discover, by considering the period of the inscriptions to either side of it, how
Shima believed it should be dated. An example is provided below. Second, the more recent
collections of rubbings and the more recent commentaries to older collections either arrange
their rubbings chronologically or assign a date to most of the inscriptions that they gloss.
The collections organized in this way are listed in table 13. Third, if the inscription in
question (or an inscription on the same bone or shell) contains the name of a diviner, it is
possible to determine the period to which that diviner is generally thought to belong. This
information may either be found by consulting one of the dated diviner lists with which
the literature abounds such a list is provided in table 6); by consulting \pinyinterm{rao-zongyi}'s
\pinyinterm{yindai-zhenbu-renwu-tongkao} (sec. 3.3.2), which assigns periods to the diviners and presents
the reasons for doing so: or by consulting the list of diviners and their inscriptions given in
\$557-576 and then, by consulting the appropriate commentaries, dating some of the other
inscriptions made by that diviner.
Other rules of thumb are given in the body of this chapter, but these three are the
quickest and simplest to apply. For example, we wish to date \pinyinterm{yinxu-title} 16 (D) (fig. 24). Using
the procedures described in\ref{sec:chapters-ch03:3.6.1} (see  sec. 3.6.1), the inscription may be found at \$289.4 (fig. 25). The
inscription immediately to the right of \pinyinterm{yinxu-title} 16 is \pinyinterm{yinxu-title} 1535 (D) which contains the
name of the diviner Lü). S567 lists all the inscriptions made by this diviner. Selecting
those inscriptions for which we know there are glosses that assign periods, we find that Ikeda
\listitem{1964} dates \pinyinterm{houbian} 1.5.8 to period IIb; similarly, all the \pinyinfull{jingjin} inscriptions (spanning
3175 to 3754) and \pinyinfull{xucun} inscriptions (spanning 1.1478 to 1.1725; 2.599 to 2.728) listed
at S567 are assigned, by \pinyinterm{hu-houxuan}'s ordering of these collections (see table 13) to
period II. Table 6 also reveals that Lü was a period II diviner. \pinyinterm{rao-short} (1959), p. 922, lists the
ancestral titles recorded in Lü's divinations that permit us to place Lü in period IIb.
\pinyinterm{yinxu-title} 1535, therefore, may be taken as a period IIb inscription. The first five inscriptions
reproduced to the left of \pinyinterm{yinxu-title} 16 at \$289.4 (fig. 25) yield, in Shima transcription at
least, no obvious clue as to their date; the sixth inscription, however, is \pinyinterm{houbian} 1.14.4,
dated by Ikeda \listitem{1964} to period IIb. It may also be determined that (邢 ), the
diviner of this inscription, is generally dated to period II or IIb; this may be done either
by consulting table 6, by consulting \pinyinterm{rao-short} (1959), p. 1029, or by consulting some of the
glossed inscriptions listed at \$567–568. \pinyinterm{yinxu-title} 16 is thus bracketed in Sōrui by two inscriptions from period IIb. In Shima opinion, therefore-which relies, of course, on the
opinions of other scholars-it should also be dated to IIb.

\subsection[Five Periods Introduction]{\origsecnum{4.1.1}Five Periods Introduction}
\label{sec:chapters-ch04:4.1.1}

inscriptions by the reigns of individual kings or groups of kings.\footnote{In some 
cases, two or more inscriptions may record a full \pinyinterm{shang} date-cyclical day, month, 
ritual cycle number (sec. 4.3.1.7). In such cases, or when inscriptions come from 
the same bone or, although on different bones, refer to the same topic (\ref{ch:3} (see ch. 3), nn. 86, 
108), it is possible to establish their relative date with exactitude. Such determinations 
which, once they involve more than one \pinyinterm{shang} month, require us to reconstruct at least 
part of the \pinyinterm{shang} luni-solar calendar-have been proposed in great number by 董 \listitem{1945} 
(see my comment at \ref{ch:3} (see ch. 3), n. 109). In this chapter, however, I am concerned only with 
macroperiodization.} The basic periodization from which all others derive, even when they 
reject some of its features, is \pinyinterm{dong-zuobin}'s system of five periods. The periods and their 
elaborations are given in table 14. To date an inscription to period II, for example, means 
we think it was made either in the reign of King \pinyinterm{zugeng} or King \pinyinterm{zujia} but that we cannot 
determine which. When it is possible to identify a specific reign, this is usually done. Thus, 
\pinyinterm{chen-mengjia-spaced} has proposed a nine-period scheme, with one period for each of the nine kings 
from \pinyinterm{wu-ding} to \pinyinterm{di-xin}.\footnote{\pinyinterm{chen-mengjia-spaced} (1956), p. 138. Few, if any, inscriptions 
can be firmly dated to the generation prior to \pinyinterm{wu-ding} (see n. 24, but cf. too n. 150), so 
that \pinyinposs{dong-short} period I should, for most purposes, be reduced to the reign of \pinyinterm{wu-ding}. For the 
possibility that \pinyinterm{lin-xin} did not reign, so that there were only eight rather than nine kings 
in the historical period, see table 1, note h.} 

Rather than speaking of periods I through IX, however, many modern scholars have preferred 
to subdivide \pinyinposs{dong-short} five periods by adding the postfixes a and b as indicated in the second 
column in table 14; thus, an inscription made in the reign of King \pinyinterm{wen-wu-ding} is dated not 
to what would logically be period VII but to period IVa. Since these conventions are generally 
accepted, it seems best to continue using \pinyinposs{dong-short} categories, which have the advantage of 
permitting us to retreat to, for example, a more general period II or period III+IV date when 
unable to determine the period more precisely. \pinyinterm{hu-houxuan}, in fact, proposed a four-period 
scheme which, due to the difficulty of separating the inscriptions of the two periods, combines 
\pinyinposs{dong-short} periods III and IV. \pinyinterm{dong-short} himself came to adopt a different four-period system, with the 
first period being of the Old School (his old periods I and IIa), the second being of the New 
School (IIb and III), the third of the Old (IV), and the fourth of the New (V). In cases for 
which the five periods themselves are still too precise, \pinyinterm{chen-mengjia-spaced}, following certain 
trends in the character of the inscriptions, has proposed a broader periodization: Early 
(periods I through IIIa), Middle...

[Content continues with full footnotes]
\section[\pinyinterm{dong-zuobin}'s Five Periods]{\origsecnum{4.2}\pinyinterm{dong-zuobin}'s Five Periods}
\label{sec:chapters-ch04:4.2}

Since most of the inscriptions cannot be placed with exactitude on a calendar of
years, scholars have proposed broader schemes of periodization which attempt to date the
% source: scan 110, printed 90
inscriptions by the reigns of individual kings or groups of kings.3 The basic periodization
from which all others derive, even when they reject some of its features, is \pinyinterm{dong-zuobin}'s
system of five periods. The periods and their elaborations are given in table 14. To date an
inscription to period II, for example, means we think it was made either in the reign of King
\pinyinterm{zugeng} or King \pinyinterm{zujia} but that we cannot determine which. When it is possible to
identify a specific reign, this is usually done. Thus, \pinyinterm{chen-mengjia-spaced} has proposed a
nine-period scheme, with one period for each of the nine kings from \pinyinterm{wu-ding} to \pinyinterm{di-xin}.4
Rather than speaking of periods I through IX, however, many modern scholars have
preferred to subdivide \pinyinposs{dong-short} five periods by adding the postfixes a and b as indicated in
the second column in table 14; thus, an inscription made in the reign of King \pinyinterm{wen-wu-ding}
is dated not to what would logically be period VII but to period IVa. Since these conventions are generally accepted, it seems best to continue using \pinyinposs{dong-short} categories, which have
the advantage of permitting us to retreat to, for example, a more general period II or
period III+IV date when unable to determine the period more precisely. \pinyinterm{hu-houxuan},
in fact, proposed a four-period scheme which, due to the difficulty of separating the
inscriptions of the two periods, combines \pinyinposs{dong-short} periods III and IV. \pinyinterm{dong-short} himself came
to adopt a different four-period system, with the first period being of the Old School (his
old periods I and IIa), the second being of the New School (IIb and III), the third of the
Old (IV), and the fourth of the New (V). In cases for which the five periods themselves
are still too precise, \pinyinterm{chen-mengjia-spaced}, following certain trends in the character of the
inscriptions, has proposed a broader periodization: Early (periods I through IIIa), Middle
C C
C C C
LCㄈㄈ
3. In some cases, two or more inscriptions may
record a full \pinyinterm{shang} date-cyclical day, month,
ritual cycle number (sec. 4.3.1.7). In such cases,
or when inscriptions come from the same bone or,
although on different bones, refer to the same
topic (\ref{ch:3} (see ch. 3), nn. 86, 108), it is possible to establish
their relative date with exactitude. Such determinations which, once they involve more than
one \pinyinterm{shang} month, require us to reconstruct at
least part of the \pinyinterm{shang} luni-solar calendar-have
been proposed in great number by 董 \listitem{1945}
(see my comment at \ref{ch:3} (see ch. 3), n. 109). In this chapter,
however, I am concerned only with macroperiodization.
4. \pinyinterm{chen-mengjia-spaced} (1956), p. 138. Few, if any,
inscriptions can be firmly dated to the generation
prior to \pinyinterm{wu-ding} (see n. 24, but cf. too n. 150),
so that \pinyinposs{dong-short} period I should, for most purposes,
be reduced to the reign of \pinyinterm{wu-ding}. For the
possibility that \pinyinterm{lin-xin} did not reign, so that
there were only eight rather than nine kings in
the historical period, see table 1, note h.
``period Ib'' for the RFG inscriptions (\ref{sec:chapters-ch02:2.2.1.1} (see see sec.
2.2.1.1)); Zhang Zongdong (1970), p. 21, follows
this
pattern. \pinyinterm{xu-jinxiong} (1974), p. 211,
introduces further subdivisions, such as IIIa,
IIIb, and IIIc.
for
6. \pinyinterm{jingjin}, preface, p. 1b; \pinyinterm{xucun}, preface,
p. 2. On the need for combining periods III and
IV, see too \pinyinterm{renwen} shakubun, English preface, p. 21;
\pinyinterm{yan-yiping} (1961), pp. 485-486. \pinyinterm{chen-mengjia-spaced} (1956), p. 139, rejects \pinyinterm{hu-houxuan}'s
solution, but not, in my view, effectively. On the
difficulties involved in separating III and IV,
consider Shima apology ([1960], p. 49
assigning certain inscriptions to period III in his
book, which he would now place in IVb@@
Similarly, the assignment of a period IV date
to certain \pinyinterm{jiabian} and \pinyinterm{renwen} inscriptions has
been challenged, plausibly, in my opinion see
n. 44). For a new attempt to distinguish periods[\textasciicircum{}5]
III and IV on the basis of physical criteria, such
as hollow shapes and burn marks, see nn. 155 and
170 below.

[\textasciicircum{}5]: This is the system adopted by Japanese 7. \pinyinterm{dong-short} (1945), pt. 1, \ref{ch:1} (see ch. 1), p. 2b; pt. 2, scholars in particular. Ikeda \listitem{1964} reserves p. 2b. Cf. \ref{ch:2} (see ch. 2), n. 18, above. \ref{ch:5} (see ch. 5), U U J
% source: scan 111, printed 91
(IIIb through IVb), and Late (V).8 These dating\footnote[8]{\pinyinterm{chen-mengjia} (1956), pp. 138, 142.}\footnote[9]{Foreign statelets and persons, for example, may be considered under divination topic (sec. 4.3.1.12; appendix 5, secs. 3 and 4); genealogy (by which 董 means \pinyinterm{fugeng} ) and titles (by which he means \pinyinterm{pangeng}) may be considered together under ancestral titles (sec. 4.3.1.1).}\footnote[10]{Courageous editors, such as \pinyinterm{hu-houxuan} (\pinyinterm{jingjin}, \pinyinterm{xucun}, and 刘-lu), Kaizuka and Itō (\pinyinterm{renwen}), or \pinyinterm{xu-jinxiong} (Menzies and White), who seek to date, often by analogy, every inscription in their collections, risk a certain number of errors. (On this point, cf. Serruys [1974], p. 15.) \pinyinterm{xu-jinxiong}, for example, originally placed Menzies 2622 in period IV, but subsequently ([1973], pp. 25, 122) assigned it to period III; he originally assigned Menzies 12 to period but subsequently joined it to Menzies 2995 of period V (\pinyinterm{zhuixin} 549.}\footnote[11]{The alternation of the Old and New Schools proposed by \pinyinterm{dong-short} (see \ref{ch:2} (see ch. 2), n. 18) provides such a theory, but it is not, in my view, satisfactory, in part because it assumes, for no demonstrable reason, that the periodizing criteria were strictly related to royal reigns. As \pinyinterm{li-xueqin} (1957), p. 124, has noted, the oracle-bone inscriptions of a single king are not necessarily all in the same style; and inscriptions in the same style do not necessarily belong to the same king.}
% source: scan 112, printed 92

\subsection*{Ancestral Titles}

The kin titles by which the \pinyinterm{shang} inscriptions refer to the dead comprise the only
unequivocal criterion for relative dating we so far possess.\footnote[12]{The best introduction to the use of ancestral and kin titles is \pinyinterm{chen-mengjia-spaced} (1956), pp. 367-501. It is to \pinyinterm{wang-guowei} ([1923], ch. 9) that the great credit goes for first using the ancestral titles to periodize the inscriptions; his study was first published in 1917 (for the bibliographical details, see \pinyinterm{hu-houxuan} [1952], 426).} 12 That is, a title, if it is
unambiguous in its referent, is prescriptive; it permits us to establish with absolute certainty
the generation when a divination, and presumably its record, was made. All other criteria,
such as diviners' names, epigraphic forms, hollow shapes, or pit location, are-to the
extent that their conclusions are expressed in terms of \pinyinposs{dong-short} five periods-secondary;
they finally derive from the evidence of the titles and, being descriptive, permit us to
establish the period of inscriptions without datable titles only with varying, but frequently
high, degrees of probability.
The way in which we can use titles to date inscriptions is simple: if an inscription
refers to an ancestor as ``father'' (fu¼) and we can place that ``father'' in the \pinyinterm{shang} royal
genealogy (table 1), we can then establish beyond doubt that the inscription refers to a
divination performed during the generation of his ``sons.'' Similarly, reference to an ``older
brother'' (hsiung ) establishes that the divination was performed by a ``younger brother”
of the same generation. The same considerations apply to titles like ``mother'' (mu),
``grandfather'' (tsu), “grandmother” or, occasionally, ``mother'' (pi), and ``son''
(tzu). These kinship terms are placed in quotation marks to indicate that they were
classificatory, not just biological.
To locate such generational titles on the royal family tree, thus establishing the period
of the inscriptions in which they occur, is, however, less easy than it may sound. The
difficulties derive partly from our inability to identify with certainty those lesser branches
and twigs of the \pinyinterm{shang} genealogy that are recorded in the inscriptions but not in 史赤¸¹³
And they derive more seriously from the ambiguity of the kin titles used in the inscriptions.
Tsu and pi are by their nature ambiguous since they were used to refer to any grandfather
and grandmother two or more generations distant from the living king; thus, both \pinyinterm{wu-ding}
(K21)* in period I and \pinyinterm{di-xin} (K28) in period Vb referred to K13 as \pinyinterm{zu-xin}, [\textasciicircum{}13]“Grand-

[\textasciicircum{}13]: The names of the dead \pinyinterm{shang} kings recorded in the oracle inscriptions have, in most cases, been equated to those given by \pinyinterm{sima-qian} in 史赤, \ref{ch:3} (see ch. 3), "Yin pen-chi"; the generational relationships recorded in 史赤, are confirmed by the kinship and sacrificial titles employed in the inscriptions. The question of whether the \pinyinterm{shang} genealogy, especially its more remote sections, was a true king list, as \pinyinterm{sima-qian} regarded it, or a retrospective invention of the late \pinyinterm{shang} kings (cf. Kane [1973], p. 357) has no bearing on its use as a dating criterion for the historical period. For a detailed comparison of the genealogical data contained in the inscriptions and in 史赤, see the notes to tables I and 15. *The K numbers refer to the kings' accession order. K21, for example, refers to the twenty-first \pinyinterm{shang} king after and including the dynasty founder, known as 唐 or Ta 乙 (K1). The names of the kings and their K numbers are given in table 1. The circled numbers indicate the order by which dead kings received sacrifice (see n. 21). 95 ㄈㄈㄈ□ L
% source: scan 113, printed 93
father 新''; he was \pinyinposs{wu-ding} great-grandfather, but he was \pinyinposs{di-xin} great-great-great-great-great-great-great-grandfather.\footnote[14]{For period I inscriptions, see \$527.1; for period V inscriptions, see S528.3.} 14 The titles tsu and pi, therefore, cannot be used
to date an inscription with certainty. The same ambiguity plagues even the generationally
specific titles like father and mother. For example, Fu Ting was a period II appellation
for \pinyinterm{wu-ding} (K21), a period IVa appellation for \pinyinterm{keng-ding} (K24), and a period V
appellation for \pinyinterm{wen-wu-ding} (K\footnote[15]{For period I inscriptions, see S547-4-548.3; for period IVa, S548.3-549.2; for period V, \$549.2.}26).15 \pinyinterm{fuyi} was a period I appellation for \pinyinterm{xiaoyi}
(K20) and a period IVb appellation for \pinyinterm{wu-yi} (K\footnote[16]{For period I inscriptions, see S545.4-547.3; for period IVb, S547.3. Note that many of these inscriptions such as \pinyinterm{jiabian} 2046; \pinyinterm{qianbian} 6.19.5, \pinyinterm{renwen} 3013; 3014; 3015; 3017-dated by Shima placement in Sōrui to period IV, are RFG divinations which would be dated by many scholars to period I or Ib (see \ref{ch:2} (see ch. 2), n. 18), with \pinyinterm{fuyi} again referring to \pinyinterm{xiaoyi} (K21); cf. \pinyinterm{yan-yiping} (1961), pp. 516, 519.}25).16 Scholars have provided lists of
the various kin titles that appear in each period,\footnote[17]{\pinyinterm{dong-short} (1933), facing p. 344; (1965), facing p. 78; Kaizuka and Itō (1953), facing p. 78; \pinyinterm{chen-mengjia-spaced} (1956), p. 156; \pinyinterm{renwen} shakubun, English preface, p. 17. Cf. S555, 556.} 17 but the generational recurrence of
titles such as \pinyinterm{fugeng} and Xiong Ting prevents these lists-which in any event may
reflect the author's own interpretations of several controversial cases-from serving as
automatic touchstones for periodization. There are few. if any, kinship terms that by
themselves may be used to date an inscription unequivocally. Each kin term and the
context in which it appears must be considered individually.
The ambiguity of the individual titles requires us to analyze title clusters. And our
analysis depends upon the assumption that the \pinyinterm{shang} worshipped their ancestors in groups
that were selected according to one or more rational principles-rational, that is, in terms
of \pinyinterm{shang} social structure and religious belief. Two principles may be discerned:
I. Vertical. Some divinations group the ancestors of certain segments of the ta-tsung
, or main direct lineage, in their appropriate order.\footnote[18]{To belong to the ta-tsung, a dead king had to have both a father and a son who had been or were kings; the predynastic ancestors of the main lineage were also regarded as members of the ta-tsung. For an introduction to the ta-tsung system, see \pinyinterm{chen-mengjia-spaced} (1956), pp. 460-468; \pinyinterm{zheng-dekun} (1960), p. 221; \pinyinterm{huang-ranwei} (1967), pp. 89-119; Ho (1975), p. 292.} 18 For example, the charge
``On
the next yi-yu (day 22), we will offer yu and fa sacrifice to (the spirits of)\footnote[19]{For ``ancestral spirits'' as a translation of see \ref{ch:1} (see ch. 1), n. 71. For a series of rituals offered to ta-tsung ancestors, see, for example, \pinyinterm{cuibian} 176 (fig. 18).} 19 the five ancestral
altars: \pinyinterm{shang} 甲, 成, Ta Ting, Ta 甲, Tsu 乙\footnote[20]{\pinyinterm{pinbian} 41.19.}'' 20 lists the ancestors, all of the main
lineage, in the order: C:①:D:D:DB: KI⑦; K2 ⑧; K3 ⑨\%; K\footnote[21]{The circled numbers indicate the ritual order in which ancestors (not all of whom had necessarily come to the throne) received sacrifice in the sacrificial schedules of periods IIb and V; on the sacrifice schedule, see \pinyinterm{chen-mengjia} (1956), pp. 386-392. These circled numbers, together with their corresponding K numbers (see the note at the bottom of p. 95), are given in tables I and 15.}12 ①8.21 In such cases, generational
order was normally maintained, the ancestors being listed in order of decreasing seniority.\footnote[22]{E.g., 甲图’u 87; \pinyinterm{pinbian} 204.21; 追 15 (\$516.3, but wrongly cited as \pinyinterm{hebian} 15); \pinyinterm{houbian} 1.20.5 (\$526.1); \pinyinterm{cuibian} 112 (S516.3); \pinyinterm{yicun} 917 (S516.3). There are exceptions, in which ancestors were listed in order of increasing, rather than decreasing, seniority; e.g., \pinyinterm{kufang} 1206 (D); 1316 (D); \pinyinterm{nanbei}, ``Ming'' 352 (D) (all S341.3); \pinyinterm{yicun} 536 (\$524.2); I have not discovered an explanation for this reversal of what I take to be the more normal order.} 22
14. For period I inscriptions, see \$527.1; for
period V inscriptions, see S528.3.
15. For period I inscriptions, see S547-4-548.3;
for period IVa, S548.3-549.2; for period V,
\$549.2.
16. For period I inscriptions, see S545.4-547.3;
for period IVb, S547.3. Note that many of these
inscriptions such as \pinyinterm{jiabian} 2046; \pinyinterm{qianbian}
6.19.5, \pinyinterm{renwen} 3013; 3014; 3015; 3017-dated by
Shima placement in Sōrui to period IV, are
RFG divinations which would be dated by many
scholars to period I or Ib (see \ref{ch:2} (see ch. 2), n. 18), with
\pinyinterm{fuyi} again referring to \pinyinterm{xiaoyi} (K21); cf. \pinyinterm{yan-yiping} (1961), pp. 516, 519.
17. \pinyinterm{dong-short} (1933), facing p. 344; (1965), facing
p. 78; Kaizuka and Itō (1953), facing p. 78;
\pinyinterm{chen-mengjia-spaced} (1956), p. 156; \pinyinterm{renwen} shakubun,
English preface, p. 17. Cf. S555, 556.
18. To belong to the ta-tsung, a dead king had
to have both a father and a son who had been or
were kings; the predynastic ancestors of the main
lineage were also regarded as members of the
ta-tsung. For an introduction to the ta-tsung system,
see \pinyinterm{chen-mengjia-spaced} (1956), pp. 460-468; \pinyinterm{zheng-dekun}
(1960), p. 221; \pinyinterm{huang-ranwei} (1967),
pp. 89-119; Ho (1975), p. 292.
19. For ``ancestral spirits'' as a translation of
see \ref{ch:1} (see ch. 1), n. 71. For a series of rituals offered to
ta-tsung ancestors, see, for example, \pinyinterm{cuibian} 176
(fig. 18).
20. \pinyinterm{pinbian} 41.19.
21. The circled numbers indicate the ritual
order in which ancestors (not all of whom had
necessarily come to the throne) received sacrifice
in the sacrificial schedules of periods IIb and V;
on the sacrifice schedule, see \pinyinterm{chen-mengjia-spaced}
(1956), pp. 386-392. These circled numbers,
together with their corresponding K numbers (see
the note at the bottom of p. 95), are given in
tables I and 15.
22. E.g., 甲图’u 87; \pinyinterm{pinbian} 204.21; 追
(\$516.3, but wrongly cited as \pinyinterm{hebian} 15);
\pinyinterm{houbian} 1.20.5 (\$526.1); \pinyinterm{cuibian} 112 (S516.3);
\pinyinterm{yicun} 917 (S516.3). There are exceptions, in
which ancestors were listed in order of increasing,
rather than decreasing, seniority; e.g., \pinyinterm{kufang}
(D); 1316 (D); \pinyinterm{nanbei}, ``Ming'' 352 (D)@@
(all S341.3); \pinyinterm{yicun} 536 (\$524.2); I have not
discovered an explanation for this reversal of
what I take to be the more normal order.
% source: scan 114, printed 94
Horizontal. Some divinations group ancestors of the same generation; \pinyinterm{yibian} 7767,
to be discussed in a moment, is a good example. In these cases too the ancestors were listed
according to the order in which they had acceded to the throne.
These principles were not followed invariably, but they represent what we conceive
to have been the rational norm. If we find, for example, reference to 祥(?) 甲,
\pinyinterm{fugeng}, and \pinyinterm{fuxin} in the same inscription (as in \pinyinterm{yibian} 7767), we assume that they
are linked by at least one of these principles. In this case, the horizontal principle seems
most appropriate; it is neater to assume that the ancestors belong to the same generation
and are thus presumably \pinyinterm{sima-qian}'s \pinyinterm{yangjia} (K17), \pinyinterm{pangeng} (K18), and
\pinyinterm{xiaoxin} (K19 rather than 台 甲 (K3) (or Wu 甲 [K14]), \pinyinterm{pangeng} (K18),
and \pinyinterm{xiaoxin} (K19). There is no way to prove that the \pinyinterm{shang} did not on occasion
scramble their ancestors irrationally; one can imagine that they might have done so in a
desperate attempt to identify the right combination of malevolent powers that might be
plaguing them. But even if this happened, the order imposed by the specific generational
titles ``father,'' ``mother,'' ``brother,'' and ``son"-greatly reduces the chance that we
would not recognize such scrambling. That is, in this case, \pinyinterm{fugeng} and \pinyinterm{fuxin} have to
belong to the same generation; if a \pinyinterm{fugeng} were to come from the previous generation,
he would have to be called \pinyinterm{zugeng}. In short, dating of the inscriptions by title depends
upon the assumption that the system of sacrifice was genealogically rational; but the
presence of such rationality must be tested in each case.
Given the fundamental importance of the ancestral titles, a detailed example of how
they are used to date an inscription will not be out of place. A pair of complementary charges
on the right and left hypoplastrons of \pinyinterm{yibian} 7767 (S532.3) reads: \listitem{1} ``Divined: ‘(We)
should offer one bovine in yu sacrifice to 祥 (?) 甲,\footnote[23]{On the readings proposed for the graph (?) and its graphic evolution, see Shima (1958), p. 72; Takashima (1973), p. 358, n. 3.} 23 \pinyinterm{fugeng}, and \pinyinterm{fuxin}''' /
\listitem{2} ``Divined: (We) should not offer one bovine in yu sacrifice to 祥 甲, \pinyinterm{fugeng},
and \pinyinterm{fuxin}.'' None of these titles appears in the 史赤 genealogy; how, then, can we
relate them to \pinyinterm{sima-qian}'s list of \pinyinterm{shang} kings?
The ritual, cyclical names, which were transmitted to \pinyinterm{sima-qian} with considerable
accuracy (see table 15 and its notes) provide the necessary clues. The appellation 祥
甲 must be related to one of \pinyinterm{sima-qian}'s 甲 kings; \pinyinterm{fugeng}, to one of his Keng
kings; \pinyinterm{fuxin}, to one of his 新 kings. The possible candidates are listed in table 16. An
asterisk marks those kings which we can automatically exclude from consideration in this
case. It is highly unlikely, for example, that \pinyinterm{fugeng} could refer to 台 Keng (K5); any
inscription which referred to him as ``father'' would have to have been made in the reigns
of \pinyinterm{xiaojia} (K6), 台 Wu (K7), or Yung Chi (K8), eight or more generations before
the period for which we have oracle inscriptions.\footnote[24]{Few, if any, oracle inscriptions from the \pinyinterm{anyang} region predate the reign of \pinyinterm{wu-ding} (K21). Following the view that the \pinyinterm{shang} moved their capital to \pinyinterm{xiaotun} during the reign of \pinyinterm{pangeng} (on this point, see the Preface), scholars have expected to find inscriptions from the reigns of the three kings preceding \pinyinterm{wu-ding} (e.g., \pinyinterm{tongxuan}, preface, p. 2b). But this has not been confirmed by the ancestral titles. To 57 97 C C C C D C D C ] KHHD+ \& KHHD+ 1 } 24 \pinyinterm{di-xin} (K28) may be excluded in
% source: scan 115, printed 95
the same way; we have, so far, no evidence of any postconquest \pinyinterm{shang} oracle inscriptions
made by \pinyinposs{di-xin} son, \pinyinterm{wu-geng}.\footnote[25]{The proposal of Miyazaki (1970), pp. 277-278, that some of the oracle-bone inscriptions found at \pinyinterm{xiaotun} might have been divined by the rulers of the \pinyinterm{zhou-dynasty} state of 韦 is without merit; the 韦 rulers could not have referred to themselves as wang, ``king,'' a term that occurs frequently as the subject of the charge in inscriptions from all periods.} 25
We may further limit the possible candidates by assembling other inscriptions in which
they appear with other ancestors of the same cyclical name. Any 甲 ancestor recorded
in the same divination charge with 祥 甲, for example, cannot have been 祥
甲. Using the information given in table 17, we may thus exclude from the list of possible
referents in table 16 Ho 談 甲 (K11), 羌甲 (K14), and \pinyinterm{zujia} (K23), and
we are left with the possibility that 祥 甲 refers to 台 甲 (K3), \pinyinterm{xiaojia}
(K6) both already excluded as too ancient-or \pinyinterm{yangjia} (K17), the most likely
candidate.
If we attempt to apply similar logic to \pinyinterm{fugeng} and \pinyinterm{fuxin}-that is, to study the
inscriptions in which they are recorded jointly with other Keng and 新 ancestors-we
reach an impasse; there are no such inscriptions. But we may pursue these two by turning
to the 史赤 genealogy and asking the simple question: At what points in the family tree
do we find a situation where a \pinyinterm{shang} king, to whose ancestors the inscription refers, would
have been able to refer to two ``fathers,'' Keng and 新. in the generation immediately
preceding his own? The entries in the right-hand column of table 18 indicate that only in
period I (and possibly III) were there Keng and 新 ``fathers'' who could have been
worshipped jointly. (The putative \pinyinterm{fuxin} of period III, not listed in 史赤 and whose
place a divination record in the reigns of \pinyinterm{pangeng}
(K18), \pinyinterm{xiaoxin} (K19), or \pinyinterm{xiaoyi} (K20),
we would need to find an inscription (or inscriptions on the same bone or shell) which, on the
basis of calligraphy, epigraphy, diviners' names,
and other considerations, would be dated to
\pinyinposs{dong-short} period I and which made reference to
Fu Ting (K15) and \pinyinterm{fugeng} (K16). No such
inscription has yet been found. 董 and his
supporters have dated one lunar eclipse inscription
to 27 March 1373 B.C., which he would place in
the twenty-sixth year of \pinyinterm{pangeng} (K18) (\pinyinterm{pinbian} 57, k'ao-shih, pp. 90-93). Even if this absolute
date were accepted (and there are reasons for
dating the eclipse considerably later [see appendix
4, sec. 4]), the question of whether \pinyinterm{pangeng}
was on the throne at that time depends upon a
series of assumptions that are, in my view, unwarranted. Similar considerations apply to \pinyinposs{dong-short}
attempt to date certain inscriptions (e.g., \pinyinterm{jiabian}
2106, k'ao-shih, p. 265) prior to \pinyinposs{wu-ding} reign
on the basis of problematical calendrical reconstructions (cf. \ref{ch:3} (see ch. 3), n. 109). On the basis of
calligraphy, epigraphy, initial preparation, and
hollow shapes, \pinyinterm{liu-yuanlin} (1974), pp. 119-121,
plates 12.1-3, 18.1, has argued that a few scapula
fragments (e.g., 5.2.66; \pinyinterm{jiabian} 2342; 2875)
predate the reign of \pinyinterm{wu-ding}. Their form is
certainly primitive (see n. 150), but the possibility
that they represent the work of untutored hands
or the remnants of an earlier tradition cannot be
excluded; there is nothing about them which
specifically requires a pre-\pinyinterm{wu-ding} date. The
lack of oracle inscriptions from the generation of
\pinyinterm{pangeng}, \pinyinterm{xiaoxin}, and \pinyinterm{xiaoyi} suggests
either (I) that these three kings did not record
their inscriptions in the same way that \pinyinterm{wu-ding}
did, or \listitem{2} that their oracle inscriptions have not
been found, or \listitem{3} that their inscriptions contain no
recognizable dating criteria (specifically, ancestral
titles), or \listitem{4} that these three kings did not have
their capital in the \pinyinterm{anyang} region, or \listitem{5} various
combinations of the above. The matter bears on
the eleven supposedly royal tombs at \pinyinterm{xibeigang}
(长 光 [1964], pp. 46-47;
Soper [1966], p. 27; Kane [1975], pp. 103-106),
for if the three pre-\pinyinterm{wu-ding} kings did not reign in
this region they may not have been buried there.
% source: scan 116, printed 97
existence is in dispute, may, in fact, be excluded from consideration; he could not have
belonged to the triumvirate of ancestors recorded by \pinyinterm{yibian}\footnote[26]{\pinyinterm{chen-mengjia-spaced} (1956), p. 456 (cf. \pinyinterm{jiabian} 2622, k'ao-shih, p. 333), and \pinyinterm{yan-yiping} (1961), p. 528, have argued that certain period III inscriptions which record a \pinyinterm{fuxin} refer to a brother of the period II kings, \pinyinterm{zugeng} (K22) and \pinyinterm{zujia} (K23); presumably he would have died before he could come to the throne. The validity of this interpretation, which is based upon secondary criteria such as diviners' names and calligraphy, has been challenged (Shima [1958], p. 47). The possibility that on the basis of ancestral titles alone, \pinyinterm{yibian} 7767 could be dated to period III, being a divination by Kings \pinyinterm{lin-xin} (K23A) or \pinyinterm{keng-ding} (K24) to their three ``fathers,'' \pinyinterm{zugeng} (K22), \pinyinterm{zujia} (K23), and the \pinyinterm{fuxin} unknown to 史赤, is excluded by the fact that 祥 甲 and \pinyinterm{zujia} (K23) appear in the same inscriptions (table 17).} 7767.) 26 The information in
table 18 suggests, therefore, that we are dealing with a period I inscription.
To return now to 祥 甲. If he were \pinyinterm{sima-qian}'s 台 甲 (K3) or Hsiao
甲 (K6), there are no suitable Keng or 新 ``fathers'' who could have been associated
with him in any rational way. If he were \pinyinterm{yangjia} (K17), however, \pinyinterm{yibian} 7767 could
then be understood as a period I divination made for \pinyinterm{wu-ding} (K21) about sacrifice to
his three uncles or classificatory ``fathers,'' 祥 甲 (K17), \pinyinterm{fugeng} (K18), and Fu
新 (K19), who would be listed, as is normal in \pinyinterm{shang} inscriptions, in the order in which
they ascended the throne: \pinyinterm{yangjia} (K17), \pinyinterm{pangeng} (K18), and \pinyinterm{xiaoxin} (K19).
We cannot tell why, in this particular case, \pinyinterm{yangjia} (K17) was referred to as
祥 甲 rather than as ``Father” 甲, as were his two brothers, ``Father” Keng and
``Father'' 新.\footnote[27]{\pinyinterm{yan-yiping} (1961), pp. 487-488, argues that ``distant ancestors"-that is, distant in time from the diviner-had to be called by nonkin titles. Thus, a distant ancestor like Yung Chi (K8) was always referred to by his nonkin name (in the inscriptions he appears as Lü Chi ), never as Tsu Chi . Conversely, \pinyinterm{yan-yiping} has argued that a reference to K18 as \pinyinterm{pangeng} (as in \pinyinterm{yibian} 8660 [S533.1]) rather than, say, \pinyinterm{fugeng} necessarily means that the divination was made some distance in time after \pinyinterm{pangeng}'s death. If this is a true rule, then it is conceivable that \pinyinterm{yibian} 7767 was made late in period I, when 祥 甲 was felt to be more distant and his kin appellation, fu, was accordingly dispensed with. Unfortunately, \pinyinterm{yan-yiping} does not define how distant a ``distant ancestor'' must be before the rule applies; his implication ([1961], p. 517) that a great-grandfather (like \pinyinterm{xiaoyi} [K20]) was always called by a nonkin term is clearly refuted by the use of ``Tsu 乙'' to refer to \pinyinposs{wu-ding} great-great-grandfather, K12 (e.g., \pinyinterm{pinbian} 41.19: \pinyinterm{houbian} 1.28.3; \pinyinterm{yicun} 873 [all S522.4]). Similar considerations apply to the use of \pinyinterm{zu-xin} for K13 (e.g., \pinyinterm{cuibian} 250 [S527.4]). \pinyinterm{yan-yiping}'s rule, nevertheless, has some validity; beyond the oneand two-generation range of fu, ``father,'' and tsu, ``grandfather,'' the \pinyinterm{shang} had to resort to other systems of nomenclature for their ancestors or fall into some confusion. No ancestor prior to Tsu 乙 (K12) is, in fact, referred to as tsu. That the naming system showed strain as ancestors grew more distant is indicated by the fact that Tsu 乙 (K12), one, of the ancestors whose name breaks \pinyinterm{yan-yiping}'s rule, was also called Kao Tsu YiZ (Ikeda [1964], 1.3.7) and 夏 乙 FZ (\pinyinterm{pinbian} 304.21; cf. Ikeda [1964], 1.8.2; Zhang Zongdong [1970], p. 203, n. 12). Presumably the \pinyinterm{shang} were seeking ways to distinguish between Tsu 乙 (K12) and \pinyinterm{xiaoyi} (K20). By period V (e.g., \pinyinterm{houbian} 1.20.5 [\$526.1]), the latter was clearly in possession of the appellation ``Tsu 乙"-again controverting \pinyinterm{yan-yiping}'s rule of distance, for \pinyinterm{xiaoyi} would at that time have been great-great-great-great-grandfather to Ti 乙 (K27); yet from at least period IIb on (S533.2), when he was a ``closer'' ancestor, he had already been called \pinyinterm{xiaoyi}. For other discussions of the way ancestral titles changed with generations, see Qu 万 (1948), p. 223; \pinyinterm{dong-short} (1951), p. 9.}27 On occasion, in period I, he was called ``father'' in the company of his
brothers.\footnote[28]{\pinyinterm{houbian} 1.25.9 concerns the offering of single rams to the three ``fathers,'' \pinyinterm{fujia}, \pinyinterm{fugeng}, and \pinyinterm{fuxin}; the inscription has generally been dated to period I (Ikeda [1964], 1.25.9; 99 99 U C C C U 11} 28 That the three uncles did form a triumvirate is confirmed by a group of
% source: scan 117, printed 98
inscriptions, also from period I, which seek to determine whether 祥 甲, \pinyinterm{fugeng},
or \pinyinterm{fuxin} was cursing the king.\footnote[29]{\pinyinterm{pinbian} 197.15-19.} 29 That 祥 甲 should be identified with 史赤's
\pinyinterm{yangjia} is confirmed by the order in which the ancestors are listed in \pinyinterm{tongxuan} 118
(S532.3) (cf. table. 17): 祥 甲 is listed after K14 and K16 and before a king
(presumably K18) on a missing fragment, who in turn comes before K19. On the assumption
that the inscriptions were genealogically rational, it is clear that 祥 甲 must
correspond to \pinyinterm{sima-qian}'s \pinyinterm{yangjia} (K\footnote[30]{Cf. 董-tsuan 118, k'ao-shih, pp. 22b-23a. Menzies and \pinyinterm{guo-moruo-spaced} apparently arrived at this identification at about the same time (see Jung Keng [1947], p. 24).}17).30
No experienced scholar attempting to date an inscription would pursue this tedious
route. Most of the alternatives proposed above would automatically be excluded on the
basis of previous knowledge of ancestral groupings. Furthermore, other criteria, such as
diviners' names and calligraphy, could have been applied to provide a far swifter and
indeed, in this case, virtually instantaneous solution. For example, the period I court
diviners, Pin and 成, appear elsewhere on \pinyinterm{yibian} 7767; the calligraphy is of the period
I style; and the plastron was found in pit YH127, a pure period I and Royal Family group
find. But I have pursued this tedious exercise to demonstrate how ancestral titles have
been used to determine in the first place that Pin and 成 were period I diviners, that
the calligraphy of this inscription was period I calligraphy, and that YH 127 did contain
period I inscriptions. For it is only from such a determination that such secondary criteria
as diviners' names, calligraphy or pit provenance derive.

\section[Physical Criteria]{\origsecnum{4.3}Physical Criteria}
\label{sec:chapters-ch04:4.3}

\subsection*{Topics and Idioms}

Rituals and sacrifices, ancestral curses, persons and statelets, military alliances, 
and numerous other topics, idioms, and graphs recorded in the divination charges frequently 
reveal a periodicity that enables us to assign a probable date to inscriptions in which they 
appear. These changes in usage are the stuff of \pinyinterm{shang} history; the oracle-bone record is 
alive with their presence.\footnote{The historical significance of specific changes may be 
different. Changes in divination formula, calligraphy, and epigraphy presumably reflected 
changes in divining and engraving personnel. The disappearance of a particular topic from 
the divination charges may be \listitem{1} real (in which case a person or statelet, for example, 
ceased to exist), or \listitem{2} apparent. If apparent, several explanations are possible: that a 
person or statelet ceased to exist as far as the \pinyinterm{shang} state or king was concerned or-which 
is not necessarily the same thing-as far as the \pinyinterm{shang} diviners were concerned, or that the 
apparent disappearance may be due only to the accident of recovery; inscriptions recording 
that topic may yet be found.}

Broad changes can be discerned in the content of the inscriptions. By period V, emphasis 
was focused upon divinations about the sacrificial schedule, the ten-day period and the night, 
and the hunt. Other areas of life, such as dreams, sickness, enemy attacks, requests for 
harvest, and the issuing of orders were divined far less frequently than they had been in 
period I, if at all (see table 29).

The forms of \pinyinterm{shang} divination also changed markedly between periods I and V. The trend 
toward terser, less detailed, and always optimistic prognostications and verifications, 
the reduced use of crack-notations and complementary pairs of charges, the disappearance 
of marginal notations, the virtual disappearance of diviners' names, the smaller number 
of calligraphic variants and divination formulas, the reduced size of the script---these 
phenomena were related to both the standardization of \pinyinterm{shang} divination procedures and 
the general shrinkage in the scope of the topics divined.

\subsection*{Physical Criteria}

Since preference for bone or shell, as well as the techniques for preparing and cracking 
them, changed with time, physical characteristics of a particular fragment may serve as 
clues about its relative date.

\subsection[Bone or Shell]{\origsecnum{4.3.1}Bone or Shell}
\label{sec:chapters-ch04:4.3.1}

At present, the distinction between scapulimancy and plastromancy is useful mainly 
to distinguish Neolithic from \pinyinterm{shang} oracle bones; it is of little use for making the 
finer separations which concern a scholar working with the inscriptions. Bovid scapulas 
were the material used most commonly in 龙 and early \pinyinterm{shang} pyromancy, followed 
in frequency by the scapulas of pig and sheep; deer scapulas were used rarely, and turtle 
shells more rarely still. Unprepared turtle plastrons have been found in early \pinyinterm{shang} strata 
at \pinyinterm{zhengzhou}, but bovid scapulas maintained their virtual pyromantic monopoly down to the 
period covered by the main corpus of the inscriptions excavated in the \pinyinterm{anyang} region when 
turtle shells at last came to be used in approximately equal numbers.\footnote{See \ref{ch:1} (see ch. 1), 
n. 16 and nn. 23-25. For a general summary, see Itō (1962a), pp. 251-252; \pinyinterm{liu-yuanlin} 
(1974), pp. 113-115.} As we have seen (sec. 1.2.3; see too appendix 2), the use of bone or 
shell in the historical period is not a sufficient criterion to periodize a particular fragment.

\subsubsection[Physical Criteria]{\origsecnum{4.3.1.1}Physical Criteria}
\label{sec:chapters-ch04:4.3.1.1}

It is a basic assumption that inscriptions found on the same bone or shell should be 
dated to the same period. Thus it is assumed that if one of the inscriptions on \pinyinterm{yibian} 7767 
is dated to period I on the basis of the ancestral titles, then all other inscriptions on that 
plastron may also be dated to period I and that their other characteristic features may be taken 
as secondary criteria for dating period I inscriptions.\footnote{The assumption that inscriptions 
found on the same bone or shell may be dated to the same period (cf. n. 33), reasonable in itself, 
appears to be supported by the evidence of the kan-chih dates which, on the same oracle bone, 
generally lie near one another in time. It is true that, in the absence of year numbers for most 
inscriptions, we cannot be absolutely certain that any series of inscriptions was not in fact 
started in one year, laid aside for, say, twenty years, and then finished in that was coincidentally 
the proper sequence in a twenty-first year. \pinyinterm{guo-moruo-spaced} (1972), p. 3, has suggested, in fact, that 
a scapula may have been used in one period, stored, then reused in a later period after the original 
inscriptions had been erased. Even if the suggestion, which has been questioned (秋 [1972], p. 43), 
is accepted, the situation it describes is virtually unique. As a general rule, scapulas and plastrons 
do not seem to have been used for divination over}

\pinyinterm{renwen} shakubun, English preface, p. 15; cf. its placement at \$550.1), the three ``fathers'' 
being taken as \pinyinterm{yangjia}, \pinyinterm{pangeng}, and \pinyinterm{xiaoxin}; \pinyinterm{luo-zhenyu} was the first to make this 
identification (\pinyinterm{wang-guowei} [1923], ch. 9, p. 13b). The same three ``fathers'' appear in 
Ogawa 1.3, also dated to period I. And all four ``fathers'' of that generation 甲, Keng, 新, 
and 乙---were recorded on the back of the set, \pinyinterm{pinbian} 12-21 (sec. 3.7).

Thus, \pinyinterm{yibian} 7767 records two divinations conducted by the diviner 成 and one by the diviner Pin, 
raising the strong presumption that these two diviners were active in period I. Detailed study, not 
necessary to reproduce here, indicates that this is correct: the ancestral titles, calligraphy, epigraphy, 
divination content, pit provenance, and so on associated with the divinations recorded by Pin and 成 
indicate that these two diviners were indeed contemporaries, active during the reign of \pinyinterm{wu-ding}.\footnote{For 
a list of the inscriptions which record the divinations of 成 and Pin, see S561562, 557-558. For a detailed 
study of these divinations and their implications for dating, see \pinyinterm{rao-short} (1959), pp. 241-344, 345-443. For a listing 
(incomplete) of the oracle bones on which the two diviners appear jointly, see ibid., p. 1214.}
\subsubsection[The Diviners]{\origsecnum{4.3.1.2}The Diviners}
\label{sec:chapters-ch04:4.3.1.2}

It is a basic assumption that inscriptions found on the same bone or shell should be
dated to the same period. Thus it is assumed that if one of the inscriptions on \pinyinterm{yibian} 7767
is dated to period I on the basis of the ancestral titles, then all other inscriptions on that
plastron may also be dated to period I and that their other characteristic features may
be taken as secondary criteria for dating period I inscriptions.\footnote[31]{The assumption that inscriptions found on the same bone or shell may be dated to the same period (cf. n. 33), reasonable in itself, appears to be supported by the evidence of the kan-chih dates which, on the same oracle bone, generally lie near one another in time. It is true that, in the absence of year numbers for most inscriptions, we cannot be absolutely certain that any series of inscriptions was not in fact started in one year, laid aside for, say, twenty years, and then finished in hat was coincidentally the proper sequence in a twenty-first year. \pinyinterm{guo-moruo-spaced} (1972), p. 3, has suggested, in fact, that a scapula may have been used in one period, stored, then reused in a later period after the original inscriptions had been erased. Even if the suggestion, which has been questioned (秋 [1972], p. 43), is accepted, the situation it describes is virtually unique. As a general rule, scapulas and plastrons do not seem to have been used for divination over}31
\pinyinterm{renwen} shakubun, English preface, p. 15; cf. its
placement at \$550.1), the three ``fathers'' being
taken as \pinyinterm{yangjia}, \pinyinterm{pangeng}, and \pinyinterm{xiaoxin};
\pinyinterm{luo-zhenyu} was the first to make this identification
(\pinyinterm{wang-guowei} [1923], ch. 9, p. 13b). The same
three ``fathers'' appear in Ogawa 1.3, also dated
to period I. And all four ``fathers'' of that
generation 甲, Keng, 新, and 乙-were recorded
on the back of the set, \pinyinterm{pinbian} 12-21 (sec. 3.7).
% source: scan 118, printed 99
Thus, \pinyinterm{yibian} 7767 records two divinations conducted by the diviner 成 and
one by the diviner Pin, raising the strong presumption that these two diviners were active
in period I. Detailed study, not necessary to reproduce here, indicates that this is correct:
the ancestral titles, calligraphy, epigraphy, divination content, pit provenance, and so
on associated with the divinations recorded by Pin and 成 indicate that these two
diviners were indeed contemporaries, active during the reign of \pinyinterm{wu-ding}.\footnote[32]{For a list of the inscriptions which record the divinations of 成 and Pin, see S561562, 557-558. For a detailed study of these divinations and their implications for dating, see \pinyinterm{rao-short} (1959), pp. 241-344, 345-443. For a listing (incomplete) of the oracle bones on which the two diviners appear jointly, see ibid., p. 1214. \pinyinposs{rao-short} hypothesis (ibid., p. 243) that Pin was still alive in period IIb depends upon the uncertain assumption that the Pin who was ordered to act in \pinyinterm{fuyin}, ``Jen'' 48 (S276.3) of period IIb was the same person as the diviner Pin (on this problem, see n. 35). In any event, there is no unequivocal evidence that Pin was still divining after period I. Similar considerations apply to \pinyinposs{rao-short} view that the diviners Tax and Xiong were active in periods I and II (ibid., pp. 812, 879). For the argument that 成 served into period II, see n. 39, example 4, below.}32
Bones or shells on which two or more diviners' names appear are particularly valuable
for periodization since they enable us to establish groups of diviners who apparently served
as contemporaries.\footnote[33]{Cf. n. 31. \pinyinterm{rao-short} (1959), pp. 1213-1236, lists diviners recorded on the same bone or shell; \pinyinterm{renwen} shakubun, introduction, p. 40, provides a partial supplement. On diviner groups, \ref{sec:chapters-ch02:2.2.1.1} (see see sec. 2.2.1.1).} 33 甲图'u 80, for example, one of the four plastrons upon which 董
based his 1931 study of diviners, contains the names of six diviners who are all accordingly
dated to the same period, period I. And it is reasonable to assume that other inscriptions
and other diviners which appear on other bones or shells where one or more of these six
diviners is represented as a diviner may also be dated to the same period.\footnote[34]{\pinyinterm{dong-short} (1931), esp. pp. 437-440; cf. \pinyinterm{renwen} shakubun, introduction, p. 33.} 34
Diviners' names have become a routine dating criterion since \pinyinterm{dong-zuobin} first
recognized these names for what they were in\footnote[35]{The acceptance of \pinyinposs{dong-short} thesis was not immediate. See, for example, 陳 春县 \listitem{1933} or Hopkins (1934), pp. 80-81, who suggested, respectively, that the graphs which 董 identified as diviners' names referred to sacred music or to the topic being divined. A good critical summary of \pinyinposs{dong-short} contribution ([1931]; [1933], pp. 344-350) is given by Kaizuka (1946), pp. 208209 for a Chinese translation, see \pinyinterm{dalu-zazhi}, 17.1 [1958], p. 28). \pinyinposs{rao-short} attempt to challenge \pinyinposs{dong-short} fundamental premise-that individual diviners may be associated with limited periodsrests upon the unverifiable assumption that every occurrence of a diviner's name found in a charge refers to the diviner who bore that name (cf. the discussion in n. 37). \pinyinterm{renwen} shakubun, introduction, pp. 27-41, provides a convincing critique of \pinyinposs{rao-short} thesis.} 1931.35 Apart from disagreement about the
dating of the Royal Family group diviners and a few others, oracle-bone scholars have
generally agreed on the periods in which the other diviners, especially those who appear
frequently, were active.\footnote[36]{On the RFG, see \ref{ch:2} (see ch. 2), n. 18. \pinyinposs{rao-short} attempts to place four of \pinyinposs{dong-short} period V diviners in period I is effectively refuted by \pinyinterm{renwen} shakubun, introduction, pp. 32-33. For another dispute, see n. 44. 101 U 11 C} 36 Table 6 lists the major diviners' names by period. Whether these
a period of more than about a year. For example,
甲图'u 80 records twenty-three ``ten-day week''
charges running from the tenth month of one year,
through the intercalary thirteenth month, and
into the fifth month of the next year; thus, the
plastron (which is incomplete) was used for at
least nine months. Since this was the third plastron
in a set of five, we may assume that four other
plastrons, with crack numbers 1, 2, 4, and 5, were
also used in the same order over the nine-month
period.
% source: scan 119, printed 100
diviners' names were the names of only one man or of a group of men is not clear.\footnote[37]{Most \pinyinterm{shang} names recorded in the inscriptions are thought to have been generic rather than the exclusive possession of individuals; they might refer to individuals, statelets and their leaders, places, the population of those statelets or places, and even, perhaps, the local powers worshipped there. (The basic studies are \pinyinterm{zhang-bingquan} [1967a] and Hayashi Minao [1968].) This raises the possibility that one diviner's name might have been used by several diviners from the same lineage or region who might have transmitted the post from father to son. (\pinyinterm{xu-jinxiong} [1974], pp. 154, 196-202, has argued, on the basis of hollow shapes and other secondary criteria, that there were two diviners named Tui É, one active in period I, the other in period IV.) There is little doubt that \pinyinterm{shang} titles, and probably \pinyinterm{shang} names, were inherited in this way. The question of the diviners' character and status, and the rationale for their changes in office, will be treated extensively in Studies.} 37 But the
diviner (or diviners) bearing a single name appears to have been active as a diviner in
clearly discernible periods; thus the possibility that more than one diviner bore the same
name does not seriously invalidate diviners' names as a dating criterion.\footnote[38]{For the view that diviners' names are an invalid dating criterion, see Lu 史-hsien (1961), p. 72.1; \pinyinterm{zhang-bingquan} (1967a), p. 774; \pinyinterm{xu-jinxiong} (1974), p. 203.} 38
On the other hand, it is not certain that the diviners' careers were rigidly compartmentalized by reign and period; the fact that a diviner's inscriptions which contain
touchstone ancestral titles may all be dated to the same period does not preclude the
possibility that some of his divinations without such ancestral titles may have spilled over
into an earlier or later period. The question of spill-over requires more study. Scholars,
using such secondary criteria as hollow shape and writing style, have suggested that
certain diviners practiced in periods adjacent to their main one; but no consensus has yet
been reached on the extent or significance of this phenomenon, or on how its existence may
be established unequivocally.\footnote[39]{\pinyinterm{xu-jinxiong} (1974), p. 203 (cf. his initial study [1972]) finds that diviners who served in two periods were ``quite common in all five periods,'' and that those who served only in one period were the exception. Further study may eventually confirm his conclusions, but at present, plausible though the idea is, much of the evidence for such spill-over is unsupported by ancestral titles, depends on secondary criteria such as calligraphy and hollow shapes, or applies only to a small percentage of any diviner's inscriptions. The following examples will introduce the scholarship, which cannot be pursued in an introductory work of this nature: \listitem{1} The diviner(s. Ta may have served from the end of period I to early III (\pinyinterm{xu-jinxiong} [1974], p. 91; see too \pinyinterm{chen-mengjia-spaced} [1956], p. 199; the view of 焦, cited in n. 32 above; Itō [1975], p. 84, n. 6 and p. 126, n. 7; cf. S566-567, 573). \listitem{2} The diviner 楚 may have straddled periods I and IIa (焦 [1959], pp. 845-846; \pinyinterm{xu-jinxiong} [1972], p. 4a). \listitem{3} Lü, Kou, and Ho J may have practiced in periods II and III, and possibly in period I (\pinyinterm{yan-yiping} [1961], p. 521; \pinyinterm{xu-jinxiong} [1974], pp. 92-93, 172-184, 193-196; cf. Zhang Zongdong [1970], p. 22, nn. 1-4, who cites inscriptions which he believes indicate that Ho, 楚, and Ta divined in periods I and III, and that Kou divined in II and III). \listitem{4} The argument that the period I diviner 成 served into period II (e.g., 董 [1952], p. 288; \pinyinterm{yan-yiping} [1961], p. 506) depends primarily upon an uncertain eclipse date (\pinyinterm{pinbian} 59.1) and the unwarranted assumption that we know the absolute dates for \pinyinposs{wu-ding} reign (cf. \ref{ch:3} (see ch. 3), n. 109). \pinyinterm{xu-jinxiong} (1974), p. 86, however, also argues for this spill-over on the basis of hollow shapes. Kaizuka and Itō (\pinyinterm{renwen} shakubun, introduction, p. 32) cite \pinyinterm{cuibian} 281 as evidence that 成 was still divining in period II. But this is uncertain for two reasons. First, it is not certain that the diviner is 成; the graph is ambiguous and S533.2 correctly records it with a question mark. Second, they date the inscription to period II because \pinyinterm{xiaoyi} (K20) was, in period I, called \pinyinterm{fuyi}; since the inscription contains a \pinyinterm{xiaoyi}, it cannot be dated to period I. But while it is true that no other period I inscription contains the title \pinyinterm{xiaoyi} (see S533.2), this does not prove that the present inscription belongs in period II; at best we can only say that it shares that particular} 39 The fact that ancestral titles rarely, if ever, confirm this
% source: scan 120, printed 101
spill-over is not a strong argument against it, for the number of nonambiguous ancestral
titles we possess for any one diviner is relatively small.\footnote[40]{A glance at the initial pages for each diviner's entry in \pinyinterm{rao-short} \listitem{1959} will reveal how few touchstone ancestral inscriptions we possess.}\footnote[41]{\pinyinterm{dong-short} (1945), pt. 2, \ref{ch:3} (see ch. 3), pp. 26a-b, has suggested, for example, that diviners associated with divinations about the childbearing of \pinyinposs{wu-ding} consorts may have been active in the first part of period I when the king was still fertile. \pinyinterm{xu-jinxiong} (1974), p. 91, suggests that since the custom of using diviners' names was abandoned later, period III inscriptions which do record diviners' names should be placed in the first part of period III. These arguments, not conclusive in themselves, may eventually be supported by other evidence.}\footnote[42]{\pinyinterm{chen-mengjia-spaced} (1956), p. 142. \pinyinterm{xu-jinxiong} (1974), p. 9, notes that period III inscriptions which do not record the diviner's name are generally on bone rather than shell and that their calligraphy is small and angular.}\footnote[43]{\pinyinterm{xu-jinxiong} (1974), p. 9.}\footnote[44]{E.g., \pinyinterm{jiabian} 786 (S76.2); \pinyinterm{pinbian} 3: \pinyinterm{buci} 16 (S452.1); 史陶 2.185 (S364.2); \pinyinterm{yicun} 374 (S160.2); 375 (S194.3). For other examples, see table 7. It should be noted that 长 Tsungtung [1970], pp. 24-27, has argued cogently for a group of divinerless inscriptions with a pronounced style of their own which he would date not to period IV but to IIa and early IIb. He cites as examples \pinyinterm{renwen} 2258-2564 (with only one exception, not specified), which he would place not in period IV (where Kaizuka and Itō have placed them) but in period II. Zhang Zongdong's arguments may be summarized as follows: \listitem{1} The Fu Ting T 103 J}
% source: scan 121, printed 102
Calligraphy
The writing style of the inscriptions may be analyzed for convenience under two
separate headings: \listitem{1} calligraphic style, the manner in which the inscriptions were written;
\listitem{2} epigraphic structure, the shapes and component elements which formed the individual
graphs. Both criteria must be considered together when dating an inscription. The subject
is complex and deserves a monograph in its own right. What follows merely adumbrates
some of the issues and exemplifies some possible solutions.
It was another of \pinyinterm{dong-zuobin}'s major contributions to recognize that the calligraphy
of the inscriptions varied by period. He provided a series of subjective phrases to describe
the writing style of each period: that of period I, for example, was “characterized by the
predominance of large and bold characters,'' that of period II was more ``cautious and
restrained.'' His descriptions, as he acknowledged, need to be greatly refined.\footnote[45]{For the various characterizations, originally given in \pinyinterm{dong-short} (1933), pp. 421-423, see nn. 54, 55, 61, 63, and 66. \pinyinterm{dong-short} himself admitted that the use of calligraphic style to date the inscriptions was a subtle art and that his criteria were only a sketch of what was needed ([1965], pp. 100-101, 102).}45 The
evolution of oracle-bone calligraphy will not, in fact, be fully understood until a firmer,
more objective methodology involving at least the following assumptions is applied to the
data:
I. Since the diviners were not the engravers (sec. 2.9.3), we should not necessarily
expect that calligraphic (or epigraphic) style changed precisely when the reign periods
(and thus \pinyinposs{dong-short} five periods) changed. An engraver might have carved inscriptions in
periods I and II, for example, using similar writing style in both periods. It seems plausible,
and it accords with the evidence, to think that spill-over occurred among engravers as
well as diviners (sec. 4.3.1.2) and to conceive of \pinyinterm{shang} calligraphy evolving as one engraver
gave way to another, rather than changing abruptly as one king succeeded another.\footnote[46]{Cf. \pinyinterm{chen-mengjia} (1956, pp. 14, 155. For evidence of epigraphic spill-over, see the numerous cases of overlapping usage in table 26. Even if a change of kings saw a wholesale change of diviners and their staffs, it is doubtful that new and orthodox graph forms could have been established overnight without lapse or error. Further, it is unlikely that teams of engraving clerks were shuttled in and out of the palace on the accession of a new king.} 46
If, as is possible, the corps of \pinyinterm{shang} engravers was relatively small, further research may
permit us to associate inscriptions with particular engravers and to date writing style not
by periods, as we do now, but by the careers of engravers.*7
who appears in these inscriptions was clearly a
powerful manes; he could not, therefore, have been
the short-lived 康 Ting (K24) who reigned
only eight years, but must have been \pinyinterm{wu-ding}
(K21), who reigned fifty-nine years (these figures
are taken from 董 [1945]; for estimates of reign
lengths not based on traditional sources, see appendix 4, sec. 5). The genealogical sequence of
\pinyinterm{nanbei}, ``Ming'' 477 (D) (S516.3), supports such
an identification. \listitem{2} The Xiong Chi and
Xiong Keng found in these inscriptions
refer to Tsu Chi (28 in the sacrificial cycle) and
\pinyinterm{zugeng} (K22), indicating a IIa or early IIb
date. \listitem{3} Numerous personalities whom we normally associate with period I were active in these
inscriptions. \listitem{4} Many of the graph forms are
identical with those of period I. \listitem{5} The calligraphy resembles the ``strong'' calligraphy of
period I. \listitem{6} The great ``purification'' rituals (yü
) against ancestral curses, offered to the ancestors
in genealogical sequence (e.g., \pinyinterm{nanbei}, ``Ming''
[D S52.3]), correspond to period IIb@@
practice.
45. For the various characterizations, originally
given in \pinyinterm{dong-short} (1933), pp. 421-423, see nn. 54, 55,
61, 63, and 66. \pinyinterm{dong-short} himself admitted that the
use of calligraphic style to date the inscriptions
was a subtle art and that his criteria were only a
sketch of what was needed ([1965], pp. 100-101,
102).
46. Cf. \pinyinterm{chen-mengjia} (1956, pp. 14, 155.
For evidence of epigraphic spill-over, see the numerous cases of overlapping usage in table 26.
Even if a change of kings saw a wholesale change
of diviners and their staffs, it is doubtful that new
and orthodox graph forms could have been established overnight without lapse or error. Further,
it is unlikely that teams of engraving clerks were
shuttled in and out of the palace on the accession
of a new king.
47. See \ref{ch:2} (see ch. 2), nn. 104, 108.[\textasciicircum{}47 \textasciicircum{}47]: See \ref{ch:2} (see ch. 2), nn. 104, 108.
% source: scan 122, printed 103
2. Given the variety of styles found on the same shell, 48 it is clear that we are not
dealing with a single calligraphic or epigraphic tradition, but with a variety of contemporaneous traditions, the work, presumably, of a number of contemporary engravers.
This being the case, it is not always useful, or even possible, to characterize the writing
styles of one period by a single general description. It is more in accord with the evidence
to look for a series of different traditions in any one period and then to consider how they
evolved in subsequent periods.
3. Any calligraphic judgments must depend, as far as possible, upon touchstone
inscriptions which may be firmly dated to a certain period by the ancestral titles they
contain (sec. 4.3.1.1) or, failing that, by the diviner's name (sec. 4.3.1.2).49
4.
Subjective or imprecise criteria cannot always be avoided, but they should be
used as rarely as possible. We may all, when other criteria are lacking, feel that a certain
inscription belongs in period IV, say, because it looks like a period IV inscription. But
such feelings cannot be used to construct a general theory of calligraphic evolution, which
must be based upon specific, objective criteria that derive from touchstone inscriptions.
Connoisseurship is not enough.
5. Calligraphic and epigraphic considerations yield descriptive generalizations for
any period, not prescriptive rules.
The subject of calligraphic evolution is highly technical. Here I offer only a prolegomenon to introduce the reader to the major calligraphic traditions and the way they
evolved. I do so by categorizing touchstone inscriptions on the basis of four criteria:
I. The size of the graphs. To characterize the size of the graphs in an inscription,
I measure the longest vertical stroke of the character *tieng M, ``to divine,'' found in most
prefaces.50 Graphs are said to be large if that stroke is over 10 mm. in length, medium if
it is 6 to 10 mm. long, small if it is below 6 mm. in length.51
2. The skill with which the graphs were engraved. Here I use the terms ``copperplate,'' ``intermediate skill,'' and ``crude'' to distinguish levels of competence. These are
subjective judgments, but the only difficulty in applying them may be expected in separating
the intermediate-skill graphs from those engraved with greater or lesser skill. I use the term
``copperplate'' to refer to the careful, well-formed graphs found on the display inscriptions
of period I, which, for analytical purposes, I take as the model against which other graphs
may be compared. 52
3. Stroke thickness. I regard stroke thickness as normal when it accords with the size
48. Table 19 indicates the variety of styles on
\pinyinterm{pinbian} 1 (on this case, cf. \pinyinterm{li-yan-spaced} [1967], pp. 2–
3). Table 20 indicates the range of sizes found on
\pinyinterm{pinbian} 14; 18; 20; in each case, the smaller graph
is less than half the size of the larger graph on the
same plastron. Cf. \ref{ch:2} (see ch. 2), nn. 104, 108.
49. On touchstone inscriptions, see n. 1.
50. See \ref{ch:2} (see ch. 2), n. 7. If no *tieng is present,
measure the height of the particle ch'i.
51. Different graph size may not necessarily
reflect the hand of a different engraver, but may,
in certain cases on the same bone or shell, reflect
the space available when the inscription was
carved (cf. \ref{ch:2} (see ch. 2), n. 90). Preliminary research
suggests that the first divinations recorded on an
oracle bone were written in a larger hand than
those which followed (see \ref{ch:3} (see ch. 3), n. 113).
52. For display inscriptions, see \ref{ch:2} (see ch. 2), n. 9o.
% source: scan 123, printed 104
of the graph, thick when it seems unnecessarily broad. The contrast is readily apparent in
inscriptions with a *tieng over 10 mm. long (table 21). The distinction between thick and
normal strokes is harder to draw, however, as graphs grow smaller; until more precise
analyses are possible, the criterion in these cases is frequently too subjective to be useful.53
4. Page design.'' This refers to the care with which the graphs were placed on the
surface of the bone or shell. I discern two major categories: a neat or compact layout and
a less-ordered and more open layout. Subjective though this criterion is, there should be
little dispute in most cases about which term to apply.
Compared to the impressionistic language of calligraphic appreciation that has
dominated the field so far, the criteria above will appear flat and mechanical. This is their
advantage. Doubtless in need of improvement, they permit us to categorize any inscription,
with hope of general agreement, on the basis of its graph size, engraving skill, stroke thickness (when applicable), and layout. Using this approach, the general evolution of \pinyinterm{shang}
calligraphy may be described as follows.

\subsection*{Calligraphy}

Three major traditions of copperplate writing-the large, the medium, and the
small were used to record the divinations of the court diviners of period I; all the inscriptions were characterized by careful, well-formed graphs (table 19). In almost every case, the
page design of the bone or shell was neat and ordered. The large or medium graphs might
be carved with thick or thin strokes; the small graphs were, of necessity, thinly carved.54
The medium and small copperplate traditions continued to thrive in period II; the
large copperplate disappeared (table 22).55 Page design, as well as the composition of the
graphs themselves, was more compact and tight than it was for the period I inscriptions
written in small calligraphy; 56 there is an impression of more efficient use of the bone or
shell surface, of less ``white space.'
``57 At the same time, however, there was some loss of
precision in the writing of the graphs themselves; engraving appears to have become
more perfunctory as it grew more routine. 58 A small, intermediate-skill calligraphy also
53. For an initial attempt to study the grooves
objectively, see \ref{ch:2} (see ch. 2), n. 102. For subjective judgments, see n. 61 below. It is important to stress
that, as \pinyinterm{zhou-hongxiang} (1973), p. 181, has
noted, stroke thickness, as seen in rubbings, may
vary with the pressure applied by the rubbing
maker or with the kind of ink and paper used.
54. \pinyinterm{dong-short} ([1965], p. 101; cf. [1964], p. 88)
described the style of period I as large, powerful,
and extremely vigorous, ``characterized by the
predominance of large and bold characters. As a
whole the strokes are strong and powerful.'' As he
noted ([1933], p. 421), virtually all the inscriptions
in \pinyinterm{qianbian}, ch. 7, are of this large and vigorous
period I tradition; so too are \pinyinterm{jinghua} 1-8. For
examples, see figs. 8, 11-14, 16, 20. For drawings,
conveniently assembled, of other period I calligraphy on shell, see \pinyinterm{renwen} shakubun, p. 138, plate
1; p. 154. plate 2; p. 165, plate 3, etc.; for period
I calligraphy on bone, see ibid., p. 283, plate 14;
p. 297, plate 15; p. 305, plate 16, etc. Here, and
in the notes that follow, I cite the drawings only
for convenience; final reference must always be
made to the rubbings and photographs and, where
possible, to the oracle bone itself\ref{sec:chapters-ch05:5.6} (see  sec. 5.6)).
55. 董 found the style of period II to be
``cautious and restrained'' ([1965], p. 101);
``mainly medium and small in size, characters
are carefully written and have an air of stately
elegance'' ([1964], p. 88).
56. Contrast \pinyinterm{jiabian} 2869 from period II with
\pinyinterm{pinbian} 14.7 or 18.7 from period I.
57. E.g., \pinyinterm{houbian} 1.7.8; 1.7.11; \pinyinterm{renwen} 1550.
58. E.g. \pinyinterm{houbian} 1.2.3+1.2.7+1.3.4; 1.7.7+
1.7.9; 1.7.8; \pinyinterm{renwen} 1253; 1262; 1293. For an
example of period II writing style, see \pinyinterm{zhuixin}
(fig. 7)@@
% source: scan 124, printed 105
existed in period II, its layout being neat or compact in some instances, less ordered or
more open in others.\footnote[59]{This tradition was also used by the engravers of the RFG inscriptions (e.g., \pinyinterm{qianbian} 8.10.1 + \pinyinterm{jingjin} 622 [D] + \pinyinterm{renwen} 3241 [for the rejoining, see \pinyinterm{renwen} shakubun, p. 750, fig. 120] + \pinyinterm{kufang} 1259 [D]; \pinyinterm{renwen} 3127; \pinyinterm{kufang} 1988 [D]). If, as I believe, these inscriptions should be dated to period Ib, they represent another calligraphic tradition which the period II engravers drew. On the issues involved, see \ref{ch:2} (see ch. 2), n. 18. upon} 59
By period III (table 23), the copperplate tradition survived only in its small form;
page design was neat, the graphs becoming spikier and more elongated.\footnote[60]{E.g., 甲图'u 42 (fig. 26).}60 The intermediate-skill tradition existed in both its medium-size and small forms, the calligraphy looser than
that of the copperplate tradition, but the layout well ordered. In all period III traditions,
considerable white space existed around inscriptions in their orderly placement. By comparison with the small graphs of period II, the small graphs of period III appear wooden
and lacking in energy.\footnote[61]{董 found the style of period III to be decadent ([1965], p. 101), ``characterized by a general decline in the quality of writing. Characters are childish and not refined'' ([1964], p. 88). According to \pinyinterm{chen-mengjia-spaced} (1956), p. 152, the style of period IIIa continued to be careful and strict like that of period IIb, but the carving had become rough and uneven, with the head and tail of every stroke tapered, the middle being thicker. By contrast, the calligraphy of periods IIIb and IV was scattered and looser than it had been in earlier periods, or than it was to be again in period V, but the carving of period IIb was delicate and the strokes of equal width. These distinctions are subtle, and by no means evident in the rubbings of the graphs 陳 cites.}61
Due to my inability-which not all scholars share-to identify period IV inscriptions
with certainty,\footnote[62]{On this issue, cf. \ref{ch:2} (see ch. 2), n. 18; \ref{ch:4} (see ch. 4), n. 6. The number of inscriptions I can date to period IV on the basis of the ancestral titles they contain is zero. Period IVa inscriptions could presumably be dated by their references to K23A as \pinyinterm{fuxin} and K24 as Fu Ting; but these appellations, which never appear on the same oracle bone, are nonspecific: \pinyinterm{fuxin} could be a period I reference to K19, Fu Ting, a period II reference to K21 or a period Va reference to K26. Similarly, period IVb inscriptions would have to be dated by their references to K23A as \pinyinterm{zu-xin}, K24 as Tsu Ting, and K25 as \pinyinterm{fuyi}; but again, with no triple ancestral clusters on the same bone which would exclude other candidates (e.g., K13, K19, K15, K21, K20, K27) for such titles, other datings are possible. (On the ambiguity of the titles for period IV, see \pinyinterm{renwen} 1815, shakubun, p. 465, n. 1: \pinyinterm{yan-yiping} [1961], p. 515.) Most IVb datings have been made on the basis of calligraphy, and it is argued on calligraphic grounds that Fu Ting, say, must be K24. But this puts the calligraphic cart before the ancestral horse. We cannot use calligraphic style, unsupported by firm ancestral identifications, to date ancestral titles that are themselves ambiguous. \pinyinterm{yan-yiping}, in fact, has dated a great many inscriptions to period IV'a, rather than IVb, on calligraphic grounds (e.g., [1961], 141, 207-215, 218-219), presumably because these inscriptions are closer in style to the largeness and calligraphic looseness of IIIb, rather than to the smallness, neatness, and precision of V. But, lacking ancestral titles to date them positively, it is entirely possible that these inscriptions could be dated, on calligraphic grounds, to an earlier period. For example, \pinyinterm{houbian} 1.23.6 + \pinyinterm{cuibian} 9 has been dated by Ikeda \listitem{1964} to period IV, and by \pinyinterm{yan-yiping} (1961), p. 528, to IVa; but I see no firm reason for not dating the inscription, with its intermediate skill and medium-to-large-sized inscriptions, neatly arranged on the bone with some white space between them, to period III. Similarly, I would suggest that period II is as plausible as IVa for \pinyinterm{houbian} 1.4.17. Other examples of such uncertain period IV datings can readily be cited (cf. n. 44). Scholars have not hesitated to assign a period IV date to inscriptions on the basis of pit provenance, chiseled hollow shapes, absence of crack notations, or ``typical period IV calligraphy,'' but the arguments are either circular-this bone comes from a period IV pit because it has period IV calligraphy, and it has period IV calligraphy because it comes from a period IV pit-or, at best, based upon an intuitive sense of what must be a period IV inscription. I do not reject this intuitive approach. There are inscriptions which I too instinctively 107 U 11 |} 62 it is difficult to characterize the calligraphy of the period with assurance,
% source: scan 125, printed 106
but it may presumably be regarded as transitional between that of periods III and V.\footnote[63]{董 found the style of period IV to be vigorous ([1965], p. 101), a ``revival of the lively, strong style'' of period I ([1964], p. 88). He characterized ([1945), pt. 2, \ref{ch:5} (see ch. 5), p. 3b) the calligraphy of period IV as ``slender gold'' (shouchin).}\footnote[64]{E.g., small: \pinyinterm{renwen} 1795; 1804; 1815; 1816; 1829; medium-size: \pinyinterm{renwen} 1811 + \pinyinterm{cuibian} 317.}\footnote[65]{E.g., the inscriptions reproduced by \pinyinterm{renwen} shakubun, p. 506, plate 30; p. 510, fig. 80. One wonders if these changes may be related to the growing age of the period III engravers?}\footnote[66]{董 described the style of period V as strictly ordered ([1965], p. 101), the characters being ``small in size but written and carved with the minutest precision'' ([1964], p. 88).}\footnote[67]{The significance of this calligraphic evolution will be considered in Studies. For a preliminary assessment, see Keightley (1973a), pp. 48-49; (1975c), pp. 50-52.}\footnote[68]{This view, which conflicts with \pinyinposs{dong-short} belief in the alternation of the Old and New Schools, will not be accepted by all. See, for example, 乙 Kung (1957); see too \ref{ch:2} (see ch. 2), n. 18. Kaizuka (1946), p. 237, however, thinks in terms of a later coalescence of the various period I traditions.}
% source: scan 126, printed 107
that calligraphic considerations are frequently not decisive in dating a particular inscription.

\subsection*{Epigraphy}

No comprehensive study of the evolution of \pinyinterm{shang} graph forms has yet been
made.\footnote[69]{The initial study, limited and out of date, is \pinyinterm{dong-short} (1933), pp. 410-417. \pinyinterm{li-xiaoting} (1968), pp. 96-97, discusses the hypothetical evolution of the \pinyinterm{shang} graph forms. \pinyinterm{tang-lan} \listitem{1965} and \pinyinterm{li-xiaoting} \listitem{1974} provide useful introductions to the question of epigraphic evolution as a whole.}69 Table 26 presents the epigraphic changes found in forty-two graphs, most of them
frequently encountered and thus likely to be useful as secondary dating criteria.\footnote[70]{For a similar table, but covering only seven graphs, see Zhang Zongdong (1970), p. 25.} 70 The
twenty-two cyclical kan-chih graphs that appear in the preface of most inscriptions (table 7)
are particularly useful in this regard; I have included thirteen of these in table 26, excluding
the rest, which underwent little or no epigraphic change in the historical period.\footnote[71]{\pinyinposs{dong-short} table showing the evolution of the cyclical graphs ([1933], facing p. 410) has been frequently reproduced: e.g., Kaizuka (1946), p. 222; Kaizuka and Itō (1953), table 2; \pinyinterm{zheng-dekun} (1960), p. 193: \pinyinterm{dong-short} (1965), facing p. 100. It is reproduced here as table 10. As will be clear from table 26, nos. 30-42, I do not accept his conclusions in all cases.}71
In making tables of this sort, it is particularly important to avoid circular reasoning.
Many inscription dates given in commentaries depend (frequently without explicit acknowledgment) upon epigraphic rules laid down by \pinyinterm{dong-zuobin}; one cannot use such datings
to prove that \pinyinposs{dong-short} rules are correct. In preparing table 26, therefore, I have done my
best to rely on touchstone inscriptions\footnote[72]{On touchstone inscriptions, see n. 1.} 72 which may be dated by criteria unrelated to the
epigraphy of the graph in question. Table 26 can, of course, be extended virtually indefinitely. It is to be hoped that an epigraphic dictionary of oracle-bone graph forms, similar,
for example, to the work of Labat for Akkadian epigraphy,\footnote[73]{René Labat, Manuel d'épigraphie akkadienne (signes, syllabaire, idéogrammes), 4th ed. (Paris, 1963).} 73 will eventually be organized
on such chronological lines. (A few entries in \pinyinterm{jiaguwen-bian} [sec. 3.3.1] do give epigraphic
data, but these are rare). In the interim, it is always possible to chart the evolution of any
graph by working through the inscriptions provided by Sōrui (sec. 3.3.2) and checking
Shima transcriptions against the rubbings as necessary.
Some graphs grew purer and simpler with the passage of time, others more exotic
and complex.\footnote[74]{Cf. \pinyinterm{zhang-bingquan} (1952), p. 626. In some cases simplified, more symbolic graphs replaced more pictographic ones (e.g., table 26, nos. 2a, 2d; this is especially true in the case of the animal graphs (nos. 22-29) where the fuller form with the belly line was superseded by the abbreviated form without the belly line; cf. Kryukov (1973, p. 30, who notes that the graph for ``horse'' (no. 27 became more schematic as the full head was replaced by the eye and the details of mane, tail, and hooves were abbreviated (cf. n. 75. But in numerous other cases, a fuller, more baroque form, requiring more strokes, replaced or came to coexist with a simpler one (e.g., nos. 3 a, b; 5 a, b; 8 a, e: 9 a, b, d; 12 a, b; 13 a, c; 15 a, b; 16 a, c; 17a, b; 18 a, b; 19 a, b). This was particularly true in the case of the cyclical graphs (nos. 32 a, b; 34a, d; 35a, b; 36 a, b; 41 a, c; 42 a, c). The period V forms of keng, tzu, yin, ch'en, wei, shen, and yu were notably more complicated and required more strokes than the period I forms.} 74 It is hard to assess the significance of these contradictory trends,\footnote[75]{Ting Su (1966), p. 2b, has noted that the 109 U ||} 75 hard to
% source: scan 127, printed 108
discern any general principle of epigraphic change at work.\footnote[76]{\pinyinterm{dong-short} (1965), pp. 99-100, noted that late forms differed from early ones (e.g., table 26, no. 13, that elements or strokes were added or changed (e.g., nos. 15, 17, 18, 19), and that picture graphs became radical + phonetic compounds as instances of this he cites chi, which I cannot document, and feng [no. 16]). But while the evolution of individual graphs may be described in these terms, no uniform trend is discernible (cf. n. 74). \pinyinposs{dong-short} claim (1965), p. 100 (cf. [1933], p. 415; [1945], pt. 2, \ref{ch:3} (see ch. 3), p. 2b that the graphs for hsi , ``night,'' and yüeh A, ``moon'' or ``month,'' changed places in different periods is not correct (see the entries at S161.4-162.1). Graph forms (a) and (b) were both used to mean ``night'' '' (i.e., the time when the moon was visible) in all five periods, though form a was apparently becoming more prevalent by period V.} 76 It is worth noting that the
\pinyinterm{shang} engravers chose to complicate or ornament part of their script at a time when the
period V calligraphy was becoming far smaller (tables 21 and 24) and when one would
have expected them to simplify the graphs.\footnote[77]{E.g., table 26, nos. Ic; 5b; 6b; 9a; 12b; 14.1b; 17b; 18b; 19b; 22b; 25b; 27b; 31b; 32b; 34d; 35b; 36b: 39b; 41c; 42c.}
It should also be remarked that in many cases the \pinyinterm{zhou-dynasty} bronze form of a graph follows
the late \pinyinterm{shang} bone form, thus providing added support for the validity of these period V
datings. But in some cases the \pinyinterm{zhou-dynasty} bronze form derives from the \pinyinterm{shang} oracle-bone
graph used in period III+IV rather than period V.\footnote[78]{E.g., table 26, nos. 3c; 10a; 13b; 15a; 16a.}78 In many of these cases the bronze
epigraphy was conservative, following traditions which went back to the epigraphic models
of period\footnote[79]{E.g., table 26, nos. 10a; 15a; 16a.} 1.79
The general picture presented by table 26 is one of a variety of contemporary epigraphic
alternatives, some existing throughout all five periods, some for only a part of the time.
As the totals given in tables 26 and 27 indicate, the number of variants remained fairly
stable in periods I and II, rose sharply in period III+IV (when twenty-two graphs in
our sample had two or more variant forms, compared to only ten in period I), and declined
still more sharply in period V (when only five graphs had two or more variants).\footnote[80]{The trends described are misleading to the extent that they exclude the RFG inscriptions whose period is still in dispute (see \ref{ch:2} (see ch. 2), n. 18). The trends would be accentuated if the RFG inscriptions are placed in period IVb. If, on the other hand, they are contemporary with the period I inscriptions, we would have a situation in which there were approximately the same number of epigraphic variants in periods I and III + IV, with a temporary reduction in the number of variants in period II (see table 27).}80 This
suggests that the epigraphic tradition only became standardized in period V after a period
of epigraphic experiment-as fresh engravers were replacing the old?-in period III+
IV.\footnote[81]{Cf. Kaizuka and Itō (1953), pp. 66-67, who note that the greatest change in the epigraphy of the kan-chih graphs took place in period IV; thus, we may generally divide the kan-chih forms into two groups: pre-IV and post-IV. On the way in which this shift may be used to date \pinyinterm{shang} bronze inscriptions, see Kane (1973), pp. 340-342.}81 Twenty graphs had one form only in period III+IV; that number had risen to
thirty-four by period V. By the end of the historical period, the epigraphic options were
early forms of ma (table 26, no. 27) were more
pictorial than the contemporary graphs for the
other domestic animals; this suggests to him the
possibility that the \pinyinterm{shang} had the horse rather
later than they had the cow, sheep, pig, and dog;
similarly, the fact that the graphs for tiger (no. 24)
and deer (nos. 25. 26) were less simplified than
those of these four domestic animals suggests that
the \pinyinterm{shang} were relatively unfamiliar with these
wild animals. But other speculations are possible;
for example, certain animal graphs might have
been more detailed for religious reasons.
% source: scan 128, printed 109
much reduced, a fact which may reflect, as we have noted (sec. 2.9.3), a decline in the
number of engravers.
It must be stressed that epigraphic evidence (like calligraphic) can never be conclusive by itself. The presence of graph forms like numbers Ic or 7e in table 26 may suggest
a period V dating for the inscription (and in these particular cases would suggest it strongly),
but like all secondary criteria-it can never prescribe an inscription's period with
certainty. There is always the possibility that a period V form was anticipated by some
avantgarde engraver in III+IV, or that a form that is thought to have died out after periods
I and II continued to be used in III + IV as an archaistic survival, perhaps by inadvertence,
perhaps by an aged engraver set in his ways. For there is no evident reason why we should
be able to pigeonhole calligraphic or epigraphic change according to five, basically political,
periods. Such changes imply a change in the training and habits of scribes as new scribes
were introduced to replace, or work with, the old. Rather than the establishment of new
epigraphic orthodoxies at the start of a new period or periods, table 26 indicates the persistence of various epigraphic fashions, some of whose practitioners spilled over from one
period to the next, with the result that new and old styles coexisted. Many of the epigraphic
choices which eventually triumphed in period V, for example, were already in use in
III IV or earlier.\footnote[82]{On reasons for the presumed gradualness of changes in writing style, see n. 46.}89
Graphic form, in short, like the diviner's name and like calligraphic style, may provide
us with the probable floruit of a particular inscription without excluding the possibility
of a variant dating. The dating of an inscription to a certain period solely on the basis of
writing style depends on negative proof: we assign the date because no inscription with
similar calligraphy or epigraphy has yet been positively identified outside the period
assigned; but the possibility that our inscription does belong outside has not been definitely
excluded. It goes without saying that further research and the excavation of new inscriptions
may require us to revise table 26 in whole or in part.

\subsection*{Inscription Placement}

Since it reflected the divination custom of the period (being closely related to
the placement of the hollows [sec. 4.3.2.4]), as well as the decisions of the particular
engraver, inscription placement may serve as one further indication of relative date.\footnote[83]{Once again, it is important to distinguish between diviners and engravers. \pinyinterm{zhang-bingquan} (\pinyinterm{pinbian} 612, k'ao-shih, p. 95) has suggested that inscriptions which ran downward, with the columns all moving left, with no attempt at symmetry of graph placement, and with no attention paid to the central axis of the plastron are characteristic of the RFG diviner Tzu †. But in this case, as in that of \pinyinterm{yibian} 4507 (RFG), the cracks made by the diviner face each other, displaying the standard mirror symmetry. The fact that the inscriptions violate this symmetry should presumably be linked to an engraver, serving the RFG diviners as a whole-probably the engraver responsible for the similar pattern on \pinyinterm{pinbian} 611 (RFG) and not necessarily with a single diviner. III U C D} 83
No detailed study of the way the placement of the inscriptions on bone or shell evolved
% source: scan 129, printed 110
with time has yet been made; this is partly because the whole scapulas or plastrons upon
which such a study would have to depend have not been found or reconstructed in significant quantities for any period but period I. The following remarks indicate only some
of the more obvious criteria.
Inscriptions associated with the court diviners of period I and the diviners of the Royal
Family group (sec. 2.2.1.1) were frequently engraved on the back, as well as the front, of
scapulas and plastrons. This was not common in periods II to V, when inscriptions were
generally recorded only on the front.\footnote[84]{See the inscriptions cited by \pinyinterm{xu-jinxiong} 1973a), pp. 22, 110-111. See too \ref{ch:2} (see ch. 2), n. 65.}85\footnote[85]{In \pinyinterm{renwen}, the following percentages of fragments in each period have inscriptions on their back: 13 percent (period I ; 1 percent (period II; 1 percent (period III); o percent (period III + IV, period V; 0.7 percent (RFG). I am grateful to David Keegan for making these counts and those in n. 88.}
As already suggested (sec. 4.3.1.3), the page design of the inscriptions became increasingly ordered and regular with the passage of time. The order which was introduced
into the sacrificial schedule in period IIb with the advent of the New School under King
\pinyinterm{zujia} was also reflected in the more systematic placement of the inscriptions on bone
and shell.\footnote[86]{On the orderly placement of inscriptions on scapulas in periods IIb and V, see \pinyinterm{dong-short} (1945', pt. 1, \ref{ch:3} (see ch. 3), p. 11a; pt. 2, ch. 6, pp. 5b-8a. Cf. \ref{ch:2} (see ch. 2), n. 123. See too \pinyinterm{yan-yiping} (1951), p. 212; (1961, p. 539, who dates 甲图'u 209 to period IIa on the basis of the irregular position of its ten-day divination inscriptions. For the situation in period III, sce table 23.}86 Similarly, the scrambled mix of topics generally found on the period I bones
and shells well exemplified by the mixture of divinations about military campaigns,
sacrifice, and toothache found on \pinyinterm{pinbian} 12-21 (sec. 3.7)-became less common in
period II, when divinations about just one topic, such as the hunt or the sacrificial cycle,
might be clustered on a particular bone or shell over a series of days. Such pure clusters
had become the norm by period V (as in fig.\footnote[87]{\pinyinterm{hu-houxuan} (1939), pp. 432-439, gives cases, mainly from period I, of divination topics inserted between, or straddled by, others of a different nature. The custom of using a scapula or plastron for divining exclusively about the ten-day week was already present in period and continued into later periods (see the inscriptions listed at \$165.1-166.3); cf. \ref{ch:2} (see ch. 2), n. 19. For period II oracle bones used exclusively for sacrifice divinations, see \pinyinterm{zhuixin} 304 (fig. 7; translated in appendix 5, sec. 2); \pinyinterm{hebian} 22 (D: S91.3); 30 (D) (S127.3); \pinyinterm{xucun} 1.1513 (S250.2); \pinyinterm{cuibian} 176 (fig. 18). For period V oracle bones, see appendix 5, n. 7. more} 10).87
Boundary lines (sec. 2.11), common in periods I and II, had virtually ceased to be
used in period V;\footnote[88]{In \pinyinterm{renwen}, the following percentages of fragments in each period have boundary lines carved on them: 4 percent (period I); 6 percent (period II) 3 percent (period III); 0.4 percent (period III+IV); 2.3 percent (period IV); o percent (period V); 2.6 percent (RFG). If, however, following Zhang Zongdong (see n. 44), we place the period IV inscriptions in period II, consistent curve is obtained: 4.2 percent (period I); 4.5 percent (period II); 3 percent (period III); 0.4 percent (period III + IV); o percent (period V). In \pinyinterm{jiabian}, the situation is as follows: 6 percent (period I); 10.5 percent (period II); 0.8 percent (period III); 2.5 percent (period IV); o percent (period V). The figures are not fully comparable since Qu 万 does not separate the RFG inscriptions but places them in period I (on this problem, see \ref{ch:2} (see ch. 2), n. 18).}88 their disappearance may presumably be related to the increasing order
of the inscriptions which rendered the use of boundaries unnecessary.
% source: scan 130, printed 112

\subsection*{Marginal Notations}

The marginal notations carved on scapula sockets, plastron bridges, and elsewhere (sec. 1.4), which record provenance, number of pieces received, date, and the
persons responsible for the ritual preparation and for making the record, are, with some
exceptions, not found on bone or shell after period\footnote[89]{\pinyinterm{hu-houxuan} (1944), p. 67a; \pinyinterm{dong-short} (1954), p. 467. Examples appear in period III (\pinyinterm{chen-mengjia-spaced} [1956], p. 13, citing \pinyinterm{jiabian} 218, 3913.16, 3918.22; in these cases, the record had been abbreviated to the name of the recorder alone and was placed on the right xiphiplastron) and perhaps in period IV (陳, loc. cit., citing the unusual inscriptions at S363.4, whose date and meaning are not certain; see too, \pinyinterm{renwen} shakubun, pp. 355-356). For RFG marginal notations, see 陳, op. cit., p. 155.} 1.89 They were never recorded in period
V. Whether their disappearance indicates a change in the system of tribute and ritual
preparation or a change in the record-keeping method, or both, is not certain.\footnote[90]{The significance of these changes will be discussed in Studies. For a preliminary assessment, see Keightley (1973a), pp. 46-49; (1975c), pp. 4952.}90

\subsection*{Preface and Postface}

Certain variations in the formulas which preceded or followed the divination
charge (sec. 2.2) may serve as secondary dating criteria. Table 7 lists the variety of prefaces
and postfaces used, most of which were abbreviations of, or variations on, a few standard
formulas. The basic formula (table 7, no. 1.1), ``裂 on \pinyinterm{jiazi}, X (the diviner)
divined,'' appears in all periods.\footnote[91]{Some abbreviations of this formula (table 7, nos. 1.3; 2.3) also appear in all periods; others appear in periods I to III (no. 2.1.1) or I to IV (no. 3.3). The single-worded preface, *tieng, ``divined'' (no. 3.3.1), appears in all periods save that of the RFG. The ultimate abbreviationcharges recorded with no preface whatever (no. 4.3) may be found in every period. No. 1.2.2, the royal equivalent of the basic formula, appears in periods I through IV, being replaced in V by no. 1.2.}⁹1
Certain characteristic patterns are useful for periodization.\footnote[92]{E.g., formulas which included the word yüeh El ``saying,'' were limited to periods I and II (on the meaning of yüeh in such cases, see \ref{ch:2} (see ch. 2), n. 21). In this case the basic formula, common in period I, appears to have been no. 1.4: ``裂 on \pinyinterm{jiazi}, X (the diviner) divined, saying.'' The many variants or abbreviations all appear in inscriptions of period I or the RFG (table 7, nos. 1.4.2 1.4.3 1.4.4; 1.4.5; 2.4; 2.4.1\%; 2.4.3; 3.4.1; 3.4.2; 3.4.3; 3.4.4;4.2).}92 A few diviners appear
to have had predilections for certain formulas;\footnote[93]{E.g., table 7, no. 2.5.2, appears to have been a favorite preface of the period IIb diviner 邢. The RFG diviners and Shao rarely used the graph *tieng; \pinyinterm{rao-short} (1959), p. 658, notes that only 7 out of 110 inscriptions contain the *tieng and (ibid., p. 687) that only 4 of the Shao inscriptions (out of some 30 inscriptions; I take this count from S575) do so. This omission was not characteristic of other RFG diviners such as Tuié and Hsieh叶.} 93 some formulas appear to have been
restricted to the inscriptions of the Royal Family group.\footnote[94]{Table 7, nos. 1.1.1; 1.1.2; 1.2; 2.1; 2.4.1.} 94 Other formulas occur in more
than one period but tend to appear chiefly in one period.\footnote[95]{E.g., table 7, no. 1.2, appears mostly in period V; no. 1.4 appears mostly in I; no. 2.3 was common in III + IV. 113 ]]}95 Divination prefaces in periods
IV and V rarely record the diviner's name, a trend that may have been related to the
% source: scan 131, printed 113
decreasing importance of the court diviners and the increasing control of divination by
the king.\footnote[96]{See n. 90.}96
As the percentage figures at the end of table 7 indicate, the great variety of formulas
used in periods I to II was reduced by almost half in periods III to V, when many of the
earlier formulas had fallen into disuse.\footnote[97]{E.g., table 7, nos. 1.4.-1.4.6; 2.4.-2.5.1; 3.4.4.2. 98. \listitem{1} The phrase chin yüeh, ``this month,'' or chin chi yüeh, ``this nth month,'' was used only in the prefaces or postfaces of period I or the RFG; see the inscriptions at \$260.2-261.1. \listitem{2} The presence or absence of a (tsai E) in the month date may also assist periodization. Formula a. chi yüeh, ``the nth month,'' in which the month number was given without the preposition ``in'' or ``at,'' was standard in period I and in RFG inscriptions and still predominant in periods II and III+IV; it had been eliminated by period V. Formula b, tsai chi yüeh, “in the nth month,'' was rare in period I, common in II, less common in III + IV, and had become the sole and standard formula by period V (cf. \pinyinterm{chen-mengjia-spaced} [1956], p. 236; 董 [1965], p. 96). My conclusions are based upon a study of the ``in the (next) ten days there will be no disaster'' inscriptions given at S165.1-168.2. Shima listing, which is incomplete, indicates that 100 percent of the month dates used formula a in period I; that 74 percent used formula a and 36 percent used formula 6 in period II; that 81 percent used a and 19 percent used b in III + IV; and that 100 percent used b in period V. These figures are not fully accurate. For example, at least four cases of formula b appear in period I or in RFG inscriptions (\pinyinterm{qianbian} 7.43.2 [S508.4]; \pinyinterm{jingjin} 696 [D S319.2]; \pinyinterm{renwen} 3016 [\$337.4]; \pinyinterm{yibian} 15 [S467.2]) so that formula a was not the only formula in use at that time. Nevertheless, the general trend from a to b, a trend which served to lengthen and complicate the inscription, is undoubtedly real. The presence of formula a renders a period V date virtually impossible; and the presence of formula b renders a period I date unlikely. \listitem{3} The way in which the first month of the calendar year was named varied with the period. The term yi-yüeh A, ``the first month,'' was used in all periods; the term cheng-yüeh iE. A which also referred to ``the first month'' but presumably had the sense of ``the regulating or rectifying month'' that brought the start of the year into harmony with the seasons, was used only in periods II and V, and possibly in III + IV; it was especially favored by the engraver who carved the inscriptions of the period IIb diviner 邢 (cf. \pinyinterm{chen-mengjia-spaced} [1956], p. 236). For a list of the occurrences of these terms, see CKWP (1934), ``\pinyinterm{hewen},'' pp. 10b-11a; CKWP (1965), ``Howen,'' pp. 27a-b; 續\pinyinterm{jiaguwen-bian}, ``附錄,'' \ref{ch:1} (see ch. 1), p. 25. For examples by period: 乙-yüeh: period I: \pinyinterm{jiabian} 2111; \pinyinterm{pinbian} 485.14; 515.12; period II: 铁阴ün 55.1; 79.4; period IIb: \pinyinterm{jianshou} 29.5; \pinyinterm{houbian} 1.19.3; period III-IV: \pinyinterm{renwen} 1737; 1767\%; period V: \pinyinterm{houbian} 1.10.16; 2.15.1 RFG: \pinyinterm{qianbian} 6.26.1; \pinyinterm{renwen} 3242. 鄭-yüeh period IIb: \pinyinterm{kufang} 1076 (D); \pinyinterm{cuibian} 176; 1352; \pinyinterm{yicun} 401; period III+IV: \pinyinterm{wenlu} 174; 232; 723; period V: \pinyinterm{renwen} 2758; 2766. The inscriptions dated to period III + IV might be better assigned to period II (see n. 44);} 97 This indicates that by the later periods \pinyinterm{shang}
divination practice was becoming more standardized, a conclusion that accords with the
findings derived from the study of calligraphic and epigraphic evolution (secs. 4.3.1.3 and
4.3.1.4). It is interesting to note, however, that whereas period III+IV saw a peak in
the number of epigraphic variants, the peak number of prefatory formulas was reached
in period II. This suggests that divination models were standardized earlier than epigraphic
ones.
The way in which the month or ritual-cycle date was recorded in the preface or
postface may also indicate the relative date of inscriptions. In the case of the months, the
issues are technical, and full agreement has not yet been reached on which formulas belong
to which periods.98 Spill-over probably occurred, for, as with our other criteria, the variations presumably reflect the usage of different diviners or engravers who were contemporaries for a while.
% source: scan 132, printed 114
The recording of the ritual cycle itself, however, is a firmer dating criterion and
probably reflects evolving conceptions of religion and chronology. Inscriptions of the form
\inscriptionsection{INSCRIPTION}

\begin{inscriptionblock}
\inscriptionref{癸卯 divination}

(PREFACE:) 裂 on \pinyinterm{guimao} (day 40), the king divined:
(CHARGE:)
POSTFACE:
``In the (next) ten days there will be no disaster.\footnote[99]{I cannot determine whether these sentences were primarily postfaces which specified the date, i.c., it was that particular fifth month in which we performed (or were going to perform [\pinyinterm{dong-short} (1965, p. 96]) yung ritual to Ta 甲 on chiach'en or whether these sentences also served as verifications to demonstrate that the yung ritual having been performed without ``disaster fault? error? the next week had started auspiciously (cf. Keightley [1973], p. 34). To the \pinyinterm{shang}, the two functions of dating the divination and verifying its prognostication may have been commingled; even those postfaces which provide a purely secular date. e.g., ``(it was) in the eleventh month when the king attacked the Renfang--at Po'' (\pinyinterm{jingjin} 584 [D S71.2]; period V), may have implied that the successful campaign had validated the auspicious prognostication for that ten-day week.}\footnote[100]{\pinyinterm{hebian} 21 (D) (S493.2); period IIb.}\footnote[101]{Period II: e.g., \pinyinterm{houbian} 2.20.7 (\$127.4); \pinyinterm{yicun} 318 (S91.3); 906 (\$250.2). Period V: e.g., the numerous divinations at S92.1-2; 128.2; 250.4-251.1;398.4-399.1;494.3.}\footnote[102]{The five rituals were ½ (yung; S492.3 ff.), ( S250.2 ff.), 6 (chi; S91.3 ff.), (tsai S398.3 ff.), and (hsieh ; S127.2 ff.). For an introduction to the ritual system, associated with the New School diviners, see \pinyinterm{dong-short} (1945), pt. 1, \ref{ch:1} (see ch. 1), pp. 3a-3b; (1965), pp. 109113; \pinyinterm{chen-mengjia-spaced} (1956), pp. 385-392. Note that Shima \listitem{1958} argues that the cycle started with the chi ritual; \pinyinterm{xu-jinxiong} \listitem{1968} argues that it started with the yi ritual. \} 115 1}
% source: scan 133, printed 115
2月巂+素爲十世十***]<X>HODODAH •D << +D +
Xx a D \& D D C ATO
多\&王三祝\listitem{10471100}
\inscriptionsection{INSCRIPTION}
\end{inscriptionblock}


\begin{inscriptionblock}
\inscriptionref{癸未 divination}

(PREFACE:) On \pinyinterm{guiwei} (day 20), the king making cracks, divined:
(CHARGE:) ``In the (next) ten days there will be no disaster.''
(POSTFACE:) In the first month.
(PROGNOSTICATION:) The king, reading the cracks, said: ``Greatly auspicious.''
(POSTFACE:)\footnote[103]{Cf. n. 99.} 103 On \pinyinterm{jiashen} (day 21) we performed chi ritual to 祥 甲
K17 23), tsai ritual to 羌甲 (K14 20), hsieh ritual to
Qian 甲 (Kii\footnote[104]{京 [D] 518 S532.3); period V.} 17).104
\end{inscriptionblock}

Still more precise dating was provided when the postface recorded the sequential number
of the royal ritual cycle. For example:
\inscriptionsection{INSCRIPTION}

\begin{inscriptionblock}
\inscriptionref{癸巳 divination}

(PREFACE:)
(CHARGE:)
(PROGNOSTICATION:)
(POSTFACE:)\footnote[105]{Cf. n. 99.} 105
On \pinyinterm{guisi} (day 30), the king making cracks, divined:
``In the (next) ten days there will be no disaster.''
The king, reading the cracks, said: ``Auspicious.''
In the sixth month, on \pinyinterm{jiawu} (day 31), we performed yung
ritual to 江 甲 (K14); it was the king's third ritual
\end{inscriptionblock}

cycle.\footnote[106]{\pinyinterm{xubian} 1.23.5 (\$244.1); period V.}106
Inscriptions containing the phrase, wei wang chi ssu, “It was the king's nth ritual
cycle,'' have been found only in period V.\footnote[107]{\pinyinterm{chen-mengjia} 1956), pp. 233-235. In period IIb the ritual cycle, which was only 300 days long, could not yet serve as a solar calendar (Shima [1958], p. 515). It is thought that by period V the length of the cycle had expanded (with the natural increase in the number of ancestors to be worshipped as the generations accumulated) to fill 360 to 370 days, approximately the length of the solar year. The ritual cycle, therefore, may have served as a calendar, and the number of the cycle may have served to identify the reign year. The sacrificial cycle had become so regular by period V that it is possible for scholars to infer the number of the sacrificial cycle of an inscription simply from the kan-chih date and month number when a particular sacrifice was performed and thus to assign a ``year'' to the inscription, even though that information is not recorded. For examples of such reconstructions, which do not always agree, see \pinyinterm{dong-short} (1945), pt. 2, \ref{ch:2} (see ch. 2); Shima \listitem{1966} (cf. my comments to appear in RBS 12); \pinyinterm{xu-jinxiong} \listitem{1968} and (1970). The considerable problems involved will be discussed in Studies.}107
Finally, it may be noted that the place of divination was rarely recorded in period I;
prefaces (or postfaces) which record this information are generally, though not exclusively,
from period V.\footnote[108]{E.g., table 7, nos. 1.5.1; 1.5.2.}108 A period or range of periods for inscriptions which record a particular
divination site may frequently be established by consulting Sōrui (sec. 3.3.2) and the
commentaries for the date of other divinations made at the same site.\footnote[109]{E.g., the phrase tsai tui-mou pu, ``crack-making at Tui-X,'' appears only in RFG inscriptions and those of period II (RFG: Yipien 1834; \pinyinterm{qianbian} 8.6.3; \pinyinterm{xucun} 2.586 [D all S441.3]; period II: the other inscriptions at S441.2-3). By contrast, the formula pu tsai mou shih ``crack-making at X-史,'' appears mainly, if not exclusively, in period V (S442.4443.2; possible exceptions are \pinyinterm{renwen} 2141; Ninghu 1.331 [D]; 葉三 43.2). Divinations at +, at , and at, were recorded only in period V (\$417.2; 431.2-432.1; 478.4-479.2). But divinations at were recorded in both periods II and V (S472.4; period II: \pinyinterm{nanbei}, ``Ming” 395 [D]; period V: \pinyinterm{qianbian} 2.8.7). An index to important occurrences of the phrase tsai mou, ``at Xplace,'' may be found at S498; S505.3-4 lists tsai mou pu inscriptions, but is not complete.}109
% source: scan 134, printed 116

\subsection*{Prognostications}

In period I, prognostications might be auspicious or inauspicious; they might be
remarkably specific, and, in the case of the ten-day week divinations, they were invariably
inauspicious.\footnote[110]{For auspicious or inauspicious prognostications, see the inscriptions listed at \$307.1-2; for specific prognostications, see secs. 2.7, 2.8; for the invariably inauspicious, ten-day prognostications, see the inscriptions listed at S307.2–308.1; see too the discussion of display inscriptions at \ref{ch:2} (see ch. 2), n. 90.}\footnote[111]{See the inscriptions listed at \$166.4-168.2; 293.4-297.1. The existence of this change has long been recognized (e.g., Britton [1940], p. 9; Crump and Crump [1963], pp. 172–173).}\footnote[112]{E.g., 乙-chu 244 (seven charges, seven (?) identical prognostications); 乙-chu 243 (six charges, six identical prognostications); \pinyinterm{qianbian} 4.6.5 + 5.16.2, \pinyinterm{jingjin} 334 (D), and \pinyinterm{houbian} 1.19.4 (all with four charges and four identical prognostications); \pinyinterm{xucun} 1.2689 (three charges, three identical prognostications). These cases may be contrasted with \pinyinterm{xucun} 2.964 (D) (seven charges); \pinyinterm{chen-mengjia-spaced} (1956), plate 21.3 and 追 65 (both with six charges); \pinyinterm{xucun} 2.971 (D) and 乙-chu 245 (both five charges); \pinyinterm{xubian} 3.29.3, \pinyinterm{xucun} 2.966 (D) and 乙-chu 217 (all with four charges) none of these record a single prognostication. These inscriptions are all transcribed at S167.2-168.1. A count of all period V inscriptions from S166.4-168.2 indicates that of the fiftynine ten-day divination fragments which contain two or more charges, 41 percent had no prognostications at all, 41 percent had charges which were all linked to auspicious prognostications; only 18 percent were nonhomogeneous, i.e., contained charges, some of which had prognostications, some of which did not.}\footnote[113]{The lack of auspicious prognostications in most period I display inscriptions (e.g., S307.3-308.1) and the lack of inauspicious prognostications in period V indicate the divination process or divination record was being manipulated. The difference in frequency between the period I and period V prognostications is confirmed by the fact that while only 8 percent of the 1,250 period I \pinyinterm{renwen} fragments contain two prognostications on the same fragment, 54 percent of the 413 period V \pinyinterm{renwen} fragments contain two or more prognostications. That the disparity is not due to the smaller size of period V calligraphy, which would make it possible for more period V than period I inscriptions to fit on a fragment of the same size, is clear from the fact that only fifteen (7.5 percent) of the first 200 \pinyinterm{pinbian} plastrons (which are large ``fragments'' indeed) have two or more prognostications on them.}\footnote[114]{See \ref{ch:2} (see ch. 2), n. 68 and the prognostications translated in secs. 2.7 and 2.8.}\footnote[115]{E.g., \pinyinterm{qianbian} 2.35.6 (S296.2); \pinyinterm{houbian} 1.15.5 (\$296.1). }
% source: scan 135, printed 117
as the charge.116 Fifth, the change in character of the prognostications was also accompanied by a change in the graphic form of chan, ``to read the cracks.'' 9117
The situation in periods II to IV requires more study, but it seems generally true that
the prognosticatory formula, wang chan yüeh, was recorded only in inscriptions from periods
I and V (and, rarely, in inscriptions of the Royal Family group).\footnote[116]{Of the seventy-six period I prognostications on \pinyinterm{pinbian} 1-200, sixty-six (87 percent) are on the back of the shell and are separated from the charge on the front. In \pinyinterm{renwen}, twenty-one (52 percent) out of the forty period I prognostications are on the back of the bone or shell. Cf. \ref{ch:2} (see ch. 2). n. 65. In period V, by contrast, all the prognostications listed at S311.2-311.3 are recorded on the same side of the oracle bone as the charge.}\footnote[117]{Table 26, no. 7e. The of period I became the of period V (S307.3; 311.2). Commentators generally take as semantically equivalent to ✓ or its variants, , , etc. (CKWT, pp. 1099-1102).}\footnote[118]{RFG prognostications frequently omit the graph chan, e.g., ``Hsieh said: 'Auspicious'' (\pinyinterm{qianbian} 6.39.1 [S311.1]); see the other examples at \ref{ch:2} (see ch. 2), n. 69.}\footnote[119]{\pinyinterm{pinbian} 197.3 appears to be a period I verification without a prognostication, but, as in most cases of this sort, the prognostication may have been recorded on one of the other four plastrons in the set (cf. the prognostications that appear on only two of the five plastrons forming the set discussed in\ref{sec:chapters-ch03:3.7} (see  sec. 3.7)). Similar considerations may apply to various period V hunt divinations which record a verification without any prognostication; e.g., \pinyinterm{qianbian} 2.29.3 (S295.1); 2.43.3 (S294.1); \pinyinterm{jingjin} 580 (D) (S295.3).}\footnote[120]{The only exception I have noted is \pinyinterm{renwen} 1373 (period II).}\footnote[121]{E.g., \pinyinterm{pinbian} 1.3; \pinyinterm{jinghua} 2 (fig. 14); translated in\ref{sec:chapters-ch02:2.8} (see  sec. 2.8).}\footnote[122]{\pinyinterm{renwen} 1373 (S488.3); 1377 (\$488.4) (both period II); \pinyinterm{qianbian} 2.27.5 (S220.3; period V).}\footnote[123]{Verifications containing the graph (yün f), ``really, truly,'' were mainly recorded in period I (see the inscriptions listed at \$488.2491.4). RFG cases include \pinyinterm{yibian} 9067; \pinyinterm{jingjin} 2910; \pinyinterm{renwen} 3085 (all S489.1 \pinyinterm{renwen} 3091 (S489.2); \pinyinterm{qianbian} 7.8.4; \pinyinterm{nanbei}, ``Nan'' 1.60 (D) (both S489.4). Period II cases include \pinyinterm{renwen} 1373; 1374 (both S488.3); 1377 (S488.4); Houpien 2.41.12 (S489.4); 1.32.5 (S491.4); \pinyinterm{xucun} 1.1719 (S491.4); the dating is not always certain and some of the fragmentary inscriptions may be charges rather than verifications. Period III cases include \pinyinterm{renwen} 1719 (S488.3); 1922 (S488.4). \pinyinterm{renwen} 2373 (\$488.4), 2521 (S489.2), and \pinyinterm{jiabian} 620 (S489.4) have been dated to period IV, though this is not firmly established (see n. 44). Only one period V case, \pinyinterm{qianbian} 5.25.5 (\$489.1) is listed. For other period V hunt verifications, without yün, see S293.4-297.1. Any judgment}
% source: scan 136, printed 118
Attention must also be paid to the notations ziyong, “this (crack/divination
charge [?] was used (?),'' or ziyu ( = [?]). The exact meaning and function
of these terms are not understood.\footnote[124]{The basic study is \pinyinterm{hu-houxuan} (1939a), which must be supplemented by \pinyinterm{xu-jinxiong} 1963), pp. 8a-11b. It is generally agreed that ziyong and ziyu were equivalent phrases (e.g., 胡, ibid., p. 481; Ikeda [1964], 2.18.13); but it is also clear that by period V they were functionally distinct: ziyong, never used when the king was recorded as diviner, was used in sacrifice divinations (S470.3-4); ziyu was used in divinations about hunting (S41.1-42.1) and campaigning the inscriptions cited by 徐, ibid., p. 9a). Yung was certainly used as a verb of sacrifice in the inscriptions (e.g., 胡, ibid., p. 471; \pinyinterm{chen-mengjia-spaced} [1956], p. 327: \pinyinterm{jiabian} 433, k’ao-shih, p. 67). It is possible, therefore, that ziyong originally meant here (at this point in the divination process) we offered sacrifice,'' i.e., this was the charge, or this was the crack, that showed it was auspicious to proceed. Since yü, in \pinyinterm{zhou-dynasty} texts at least, meant ``to drive a chariot,'' it is conceivable that in the hunt divinations of period V ziyu meant ``here (at this point in the divination process) we drove off in our chariots'' and as a result, as the verifications frequently note, caught so many animals (徐. ibid., p. gb; cf. Ogawa 6, shakubun, p. 263, n. 5. It is hard to apply such a meaning, however, to \pinyinterm{houbian} 2.18.13, translated in this section.} 124 Since the majority of these notations, which, so far
as we can tell, always appear in auspicious contexts, were not written beside the cracks,
they cannot, strictly speaking, be classified as crack notations.\footnote[125]{That they may have come to function as crack notations is suggested by the fact that their use appears to have increased as the use of crack notations declined (sec. 4.3.1.11).} 125 By period V, at least,
the phrase ziyu served, in divinations about the weather or the hunt, as a link between
the auspicious prognostication and the verification:
\inscriptionsection{INSCRIPTION}

\begin{inscriptionblock}
\inscriptionref{壬子 hunting divination}

(PREFACE:)
On \pinyinterm{renzi} (day 49) the king, making cracks, divined:
(CHARGE:) ``Hunting at Chi, going and coming there will be no disaster.''
The king, reading the cracks, said: ``Extremely auspicious.''
Ziyu (this was used [?].'' or ``at this point we drove off in
our chariots'' [?]).¹
PROGNOSTICATION:)
(VERIFICATION\footnote[126]{See n. 124.} :)
We caught forty-one foxes, eight hornless deer, one.\footnote[127]{\pinyinterm{qianbian} 2.27.1 (S220.3; period V.} . . .
On other occasions in period V, the prognostication was omitted and the ziyu stood alone:
\end{inscriptionblock}

(PREFACE:) 裂 on \pinyinterm{guiwei} (day 20) [divined]:
(CHARGE:)
VERIFICATION:)
``This month there will be a great rain.'' Ziyu (``this was used\footnote[128]{See n. 124. ?]'').128
In the night, it rained.\footnote[129]{\pinyinterm{houbian} 2.18.13 (S41.2; period V).} 129
Whatever the exact meaning or function of these terms, which we may refer to as ``verification notations, their presence in an inscription provides a clue to its relative date. Ziyong was used in periods II to V, but has not been found in period I. Ziyu was used only
in period V.\footnote[130]{These conclusions are based upon the inscription sample at S41.1-42.2; 470.3-4. The notations tzu puyung, pu yung T, or tzu wu yung appear rarely, and, according to \pinyinterm{hu-houxuan} (1939a), p. 484, only in periods II and IV. 311 119 後下X1 88 AKESNAE} 130
about the extent to which verifications were used,
however, is complicated by various considerations: \listitem{1} The need to relate verification use to
particular divination topics, \listitem{2} the need to study
verifications not containing yün ft, \listitem{3} uncertainty
about whether the postface recording a sacrifice
date should be considered a verification (see n. 99)
-if so, such verifications also appeared, rarely, in
period II (n. 101 -\listitem{4} uncertainty about the
functions of the notations ziyong and ziyu which
may have served as verifications.
% source: scan 137, printed 119
Crack Numbers and Sets
As we have seen (secs. 2.4, 2.5), crack numbers in period I might run from numbers
I to 10. After period II, however, they never exceeded\footnote[131]{\pinyinterm{qiu-short} (1972), p. 44, notes that crack numbers never exceed number 5 in periods III and IV. And this is clearly true of period V. I have not yet found any period II sets with crack numbers higher than 5.} 5. 131 The presence of crack numbers
higher than 5, therefore, is a sign of early date.
Further, the period I practice of dividing divination topics into pairs of complementary,
balanced charges in the positive and negative mode (sec. 2.5) had virtually lapsed by
period V,\footnote[132]{There is no easy way to document this evolution since pairs are only readily identifiable on the large, reconstituted, period I plastrons such as we have in \pinyinterm{pinbian} and \pinyinterm{yibian}. We have no analogous material for later periods; more research is needed on periods II to IV, in particular. The presence of the negatives fu 弗,pu不,pi,wu勿, and wu may be taken as indication of the existence of negative charges and thus, crudely, of the existence of positive-negative charge pairs. (Wang is excluded since it frequently appears in formulaic endings, such as wang tsai, ``there will be no disaster,'' or in prognostications, but is seldom used as the negative in a complementary pair.) A count of these negatives in \pinyinterm{renwen} reveals the following situation: 221 (13 percent) of the 1,730 period I inscriptions contain negatives; 50 (19 percent) of the 269 RFG inscriptions contain negatives; 43 (10 percent) of the 431 period II inscriptions contain negatives; 201 (23 percent) of the 888 period III+IV inscriptions contain negatives; 5 (0.7 percent) of the 384 period V inscriptions contain negatives (I am grateful to Inge Dietrich and Beth Upton for making these counts). The resurgence of negatives in period III+IV is misleading here since many of the inscriptions, especially those containing the nega. tive pi-probably a substitution, used principally in period III+IV, for wu ) (Serruys [1974], p. 6; cf. \pinyinterm{chen-mengjia-spaced} [1956], p. 128-were merely single, negative charges, and not one of a positivenegative pair. This preference for the negative charge appears to have been characteristic of many period III+IV inscriptions (e.g., the inscriptions containing pi at S293.3; 380.2-383.4). And it should be noted that at least in the hunt inscriptions, the effect was to exclude certain alternatives rather than prohibit the main activity; e.g., ``Divining: ‘(The king) should not hunt on the jen day, (for if he does,) it will perhaps rain'''' (\pinyinterm{jiabian} 1920 [S293.3]), or ``(The king) should not hunt on the jen day, (for if he does,) there will perhaps be regret'' (\pinyinterm{jiabian} 3593.2 [S293.3]). By period V, in any event, complementary pairs were rarely used; I have found less than twenty cases (e.g., \pinyinterm{qianbian} 2.96.6; 2.43.5; 2.44.1; \pinyinterm{houbian} 1.20.1; 乙-chu 116 [all S457.1]; \pinyinterm{qianbian} 2.42.6; 乙-chu 115 [both S457.2]: \pinyinterm{qianbian} 2.30.6; \pinyinterm{houbian} 1.30.8 [both S455-4]; \pinyinterm{qianbian} 3.17.5 [S41.2]; \pinyinterm{cuibian} 688 [partly at \$41.2]; \pinyinterm{renwen} 2891). In most of these cases the undesirable charge (that containing the particle ch'i [see \ref{ch:3} (see ch. 3), n. 41]) was considerably abbreviated (cf. \ref{ch:2} (see ch. 2), n. 130, above). For a preliminary assessment of the significance of this change, to be considered in Studies, see Keightley (1973a), pp. 37-38,46-48,57.}132 so that period V sets consist of only five identical charges, usually couched
in terms of what the \pinyinterm{shang} wanted to have happen-``there will be no disaster,” “today
it will not rain"-with the crack numbers running from 1 to 5 (as in fig. 10). Crack numbers
and sets, in short, became standardized and simplified with the passage of time.

甲三元三I 田宮管甲九二。KIX田宮::

\subsubsection[Topics and Idioms]{\origsecnum{4.3.1.9}Topics and Idioms}
\label{sec:chapters-ch04:4.3.1.9}

Rituals and sacrifices, ancestral curses, persons and statelets, military alliances, 
and numerous other topics, idioms, and graphs recorded in the divination charges frequently 
reveal a periodicity that enables us to assign a probable date to inscriptions in which they 
appear. These changes in usage are the stuff of \pinyinterm{shang} history; the oracle-bone record is 
alive with their presence.\footnote{The historical significance of specific changes may be 
different. Changes in divination formula, calligraphy, and epigraphy presumably reflected 
changes in divining and engraving personnel. The disappearance of a particular topic from 
the divination charges may be \listitem{1} real (in which case a person or statelet, for example, 
ceased to exist), or \listitem{2} apparent. If apparent, several explanations are possible: that a 
person or statelet ceased to exist as far as the \pinyinterm{shang} state or king was concerned or-which 
is not necessarily the same thing-as far as the \pinyinterm{shang} diviners were concerned, or that the 
apparent disappearance may be due only to the accident of recovery; inscriptions recording 
that topic may yet be found.}

Broad changes can be discerned in the content of the inscriptions. By period V, emphasis 
was focused upon divinations about the sacrificial schedule, the ten-day period and the night, 
and the hunt. Other areas of life, such as dreams, sickness, enemy attacks, requests for harvest, 
and the issuing of orders were divined far less frequently than they had been in period I, if at 
all (see table 29). A full account of this evolution, which is so significant for our understanding 
of the development of late \pinyinterm{shang} religion, state, and bureaucracy, will require further research.

Specific changes in topics and idioms---changes which both exemplify, and need to be considered 
in the context of, this general evolution---are presented in appendix 5. They provide valuable 
clues for dating the inscriptions in which they appear.
\subsubsection[Crack Notations]{\origsecnum{4.3.1.11}Crack Notations}
\label{sec:chapters-ch04:4.3.1.11}

In period I, the crack notations recorded on plastrons were, listed in order of
frequency, shang-chi E, “highly auspicious,” hsiao-chi, “slightly auspicious,\footnote[133]{\pinyinterm{yibian} 7767 (left hyoplastron) contains a unique hsia-chi F, probably an engraver's error.}''
9,133 pu
% source: scan 138, printed 120
tsai ming
*135
(?), “not again charge\footnote[134]{The meaning of the crack notation S or has long puzzled commentators. My preliminary studies indicate that the notation tends to occur early in a set; for example, for the twenty-two occurrences of this notation found on \pinyinterm{pinbian} 1-512, ten are recorded by crack 1, two by crack 2, and five by crack 3, so that seventeen (77 percent) of the twenty-two cases appear by the first three cracks. Since the phrase occurs in a context of other crack notations which were auspicious (e.g., \pinyinterm{pinbian} 61.7–8; 155-3-4; 172.67; 195.1-2; 237.8; 261.1; 267.2; 345.2; 345-7; 353.1; 421.20; 491.3-4), we may assume it was either auspicious itself, or at least neutral. Since the phrase was not used after period I we are not likely to find analogues in \pinyinterm{zhou-dynasty} accounts of plastromancy. The best discussions, are those of 焦 (1957/58), pp. 16-20, and \pinyinterm{huang-peirong} (1969). Huang finds that the notation is associated with the divination inscriptions of only eight diviners, but he reaches no conclusions about its meaning. 焦 argues that \& should be read as tsai F, which is persuasive, in my view; and that should be read as ming, a loan for ming, which is more speculative; and that the entire phrase pu tsai ming meant ``not again charge'' the bone or shell because the divination had been successful. The notation, in short, may have indicated that there was no need to divine with another sequence of numbered cracks. Cf. \pinyinterm{xu-jinxiong} (1973a), pp. 48, 92, who reads it as pu wu chu 1, ``no obstacle."}\footnote[135]{This ranking is based upon a count of the first 200 \pinyinterm{pinbian} plastrons and plastron fragments, which contain 2,919 numbered cracks (based upon the information, not always exact, given in the commentaries) and 248 crack notations: 219 shang-chi notations, 17 hsiao-chi, 8 pu tsai ming (?), and 4 chi. The same order of frequency may be found on the period I \pinyinterm{renwen} bone and shell fragments, except that pu tsai ming (?) was the most commonly used notation (see n. 137); this discrepancy may be explained in part by the fact that in \pinyinterm{renwen} the notation appears twice as frequently on bone (48 instances [9 percent] on 545 fragments) as on shell (32 [4.5 percent] instances on 705 fragments); there are no bones in \pinyinterm{pinbian}. But more research, which would associate particular notations with particular topics, diviners. or engravers is needed before the relative frequency of the crack notations and its significance are fully understood.}\footnote[136]{\pinyinterm{xu-jinxiong} (1974), pp. 175-177, notes the appearance of shang-chi on \pinyinterm{xucun} 1.1603, 1.1670, and 乙-chu 393, whose diviners, Xiong, Ta, and 楚, are usually dated to period II (cf. 徐 [1973], pp. 5, 58); presumably, this is an instance of spill-over-either of period II diviners and engravers operating in period I or of period I practice continuing into period II (cf. n. 39). 徐 also argues that \pinyinterm{jiabian} 27+2 and \pinyinterm{wenlu} 78, both containing a shang-chi crack notation, but both divined by Ho, normally regarded as a period III diviner, should, on the basis of hollow shapes (sec. 4.3.2.3), be dated to late period I. üan}\footnote[137]{The basic study is 长 拼'ü (1952); in English, see Wu 史-ch'ang (1955). My preliminary conclusions are based upon a count of the crack notations listed in \pinyinterm{renwen} sakuin. In period I, I find the following notations: 80 pu tsai ming, 54 shang-chi, 20 hsiao-chi, 3 chi; on RFG fragments there is I shang-chi and I hsiao-chi; in period II, there are no crack notations; in period III+IV, I find 64 chi, 45 ta-chi, and 15 hung-chi crack notations; in period V there is one hung-chi. As with most periodic analyses of this sort, the dating of some of the inscriptions is in dispute (cf. n. 44; \pinyinterm{xu-jinxiong} [1974], pp. 13, 49, 197-198, would date the RFG fragments \pinyinterm{renwen} 3220 and 3228 to period I; cf. 徐 [1973], pp. 4, 7, 57, 68). \pinyinterm{chen-mengjia} (1956), p. 142, suggests that chi and ta-chi were especially common in period IIIb, but cites no cases; Wu 史-ch'ang (1955), p. 48, argues, by contrast, that marginal notations were rare in III but frequent in IV. More discriminating criteria may eventually permit distinctions of this sort to be made with assurance. 121 U}
% source: scan 139, printed 121
Topics and Idioms
Rituals and sacrifices, ancestral curses, persons and statelets, military alliances,
and numerous other topics, idioms, and graphs recorded in the divination charges frequently
reveal a periodicity that enables us to assign a probable date to inscriptions in which they
appear. These changes in usage are the stuff of \pinyinterm{shang} history; the oracle-bone record is
alive with their presence.\footnote[138]{The historical significance of specific changes may be different. Changes in divination formula, calligraphy, and epigraphy presumably reflected changes in divining and engraving personnel. The disappearance of a particular topic from the divination charges may be \listitem{1} real (in which case a person or statelet, for example, ceased to exist), or \listitem{2} apparent. If apparent, several explanations are possible: that a person or statelet ceased to exist as far as the \pinyinterm{shang} state or king was concerned or-which is not necessarily the same thing-as far as the \pinyinterm{shang} diviners were concerned, or that the apparent disappearance may be due only to the accident of recovery; inscriptions recording that topic may yet be found.} 138
Broad changes can be discerned in the content of the inscriptions. By period V, emphasis was focused upon divinations about the sacrificial schedule, the ten-day period and
the night, and the hunt. Other areas of life, such as dreams, sickness, enemy attacks, requests for harvest, and the issuing of orders were divined far less frequently than they had
been in period I, if at all (see table 29). A full account of this evolution, which is so significant for our understanding of the development of late \pinyinterm{shang} religion, state, and bureaucracy, will require further research.\footnote[139]{Research is needed, in particular, on the situation in periods III and IV. New School divinations about a systematic ritual schedule, which first appear in period IIb and reappear, fully developed, in period V, have not been found in significant quantities for the intervening periods. This fact, which may be due to accidents of discovery (see sec. 4.3.3.3), does not necessarily mean, however, that period IV saw a return to the Old School practices of period I. It merely means that the situation in III+IV differed from that in IIb and V. Cf. \ref{ch:2} (see ch. 2), n. 18.} 139 Specific changes in topics and idioms-changes
which both exemplify, and need to be considered in the context of, this general evolution
―are presented in appendix 5. They provide valuable clues for dating the inscriptions in
which they appear.
The forms of \pinyinterm{shang} divination also changed markedly between periods I and V. The
trend toward terser, less detailed, and always optimistic prognostications and verifications,
the reduced use of crack-notations and complementary pairs of charges, the disappearance
of marginal notations, the virtual disappearance of diviners' names, the smaller number
of calligraphic variants and divination formulas, the reduced size of the script-these
phenomena were related to both the standardization of \pinyinterm{shang} divination procedures and
the general shrinkage in the scope of the topics divined.

\subsection*{Physical Criteria}

Since preference for bone or shell, as well as the techniques for preparing and
cracking them, changed with time, physical characteristics of a particular fragment may
serve as clues about its relative date.
% source: scan 140, printed 122

\subsection[Initial Preparation Changes]{\origsecnum{4.3.2}Initial Preparation Changes}
\label{sec:chapters-ch04:4.3.2}

\subsubsection[Bone or Shell]{\origsecnum{4.3.2.1}Bone or Shell}
\label{sec:chapters-ch04:4.3.2.1}

Bone or Shell
At present, the distinction between scapulimancy and plastromancy is useful
mainly to distinguish Neolithic from \pinyinterm{shang} oracle bones; it is of little use for making the
finer separations which concern a scholar working with the inscriptions. Bovid scapulas
followed
were the material used most commonly in 龙 and early \pinyinterm{shang} pyromancy,
in frequency by the scapulas of pig and sheep; deer scapulas were used rarely, and turtle
shells more rarely still. Unprepared turtle plastrons have been found in early \pinyinterm{shang} strata
at \pinyinterm{zhengzhou}, but bovid scapulas maintained their virtual pyromantic monopoly down
to the period covered by the main corpus of the inscriptions excavated in the \pinyinterm{anyang}
region when turtle shells at last came to be used in approximately equal numbers.\footnote[140]{See \ref{ch:1} (see ch. 1), n. 16 and nn. 23-25. For a general summary, see Itō (1962a), pp. 251-252; \pinyinterm{liu-yuanlin} (1974), pp. 113-115.} 140 As
we have seen (sec. 1.2.3; see too appendix 2), the use of bone or shell in the historical period
is not a sufficient criterion to periodize a particular fragment.

\subsubsection[Initial Preparation]{\origsecnum{4.3.2.2}Initial Preparation}
\label{sec:chapters-ch04:4.3.2.2}

The initial preparation of scapulas evolved from the most primitive, Neolithic
stage, in which the bone was used in its natural state; through a transitional stage (Neolithic
and early \pinyinterm{shang}), in which the socket and spine might be cut away, frequently crudely,
and the bone might be polished; to the more developed \pinyinterm{anyang} stage, the historical period,
in which the spine was sawn away, the socket was sawn to leave the half-moon-shaped section
and the right-angled cut (sec. 1.3.1), and the surface was both smoothed and polished.\footnote[141]{\pinyinterm{liu-yuanlin} (1974), with convenient tables at pp. 110–112, 118, 119.} 141
It has been suggested, in fact, that roughly smoothed scapulas with no right-angled cut
sawn in the socket never appear after period\footnote[142]{Ibid., p. 123.} 1.142 Further research may permit relatively
subtle distinctions in preparation technique-reflecting, presumably, the habits of the
artisans involved-to be used as criteria for assigning different periods to the inscribed
bones.\footnote[143]{E.g., \pinyinterm{xu-jinxiong} (1974, pp. 209, 254-255, has proposed that variations in the shape of the ridge made on the outer edge (for the meaning of ``outer,'' see \ref{ch:1} (see ch. 1), n. 46) of scapula backs may be used to distinguish scapulas of period III from those of period IV.} 143 At present, no significant changes have been discerned in the way the initial
preparation of inscribed turtle shells (sec. 1.3.2) varied with time.

\subsubsection[Hollow Shapes]{\origsecnum{4.3.2.3}Hollow Shapes}
\label{sec:chapters-ch04:4.3.2.3}

Neolithic diviners scorched some scapulas with no advance preparation, applying the heat source directly to the unworked, natural surface; this primitive, free-form
scapulimancy continued to be practiced through at least the early \pinyinterm{shang}.\footnote[144]{Itō (1962a), p. 257, fig. 25a; \pinyinterm{liu-yuanlin} (1974), pp. 99-112. 123 == C D D ]} 144 Other Neo-
% source: scan 141, printed 123
lithic and early \pinyinterm{shang} oracle bones, however, did have pyromantic hollows, which varied
greatly in size and shape, bored into their surface, and the following trends have been
discerned in the evolution of the shapes.\footnote[145]{\pinyinterm{liu-yuanlin} (1974), p. 118, tabulates scapula hollow shapes by period, but the small size of his sample renders the general value of his conclusions uncertain.} 145 龙 scapulimancers mainly used bored
hollows that were round and shallow;\footnote[146]{Itō (1962a), pp. 251, 257; \pinyinterm{liu-yuanlin} (1974), pp. 99-112.} 146 pre-\pinyinterm{anyang} \pinyinterm{shang} diviners used hollows that
were both bored and chiseled, but with the chiseled hollow being short and small; diviners
in the reign of \pinyinterm{wu-ding} and after mainly used chiseled hollows that were long and large.\footnote[147]{\pinyinterm{liu-yuanlin} (1974), pp. 104-112, 129. Early \pinyinterm{shang} is here used to refer to oracle bones found at \pinyinterm{erligang}, 刘里-ke\pinyinterm{liulige}, 绵池, 楚-ch'iu ẞ, 奇-li-p'u , and \pinyinterm{gaohuangmiao}. Cf. the conclusions of Kaizuka and Itō who, working with the \pinyinterm{renwen} collection, advance the thesis that the bored, round hollow preceded the development of the double hollow; they further argue that cases in which the diameter of the circular hollow is equal to or greater than the vertical length of the chiseled hollow, or cases in which the long, chiseled hollow is entirely enclosed within the circular one (e.g., the back of \pinyinterm{renwen} 707; 3227; 3228; 3230; 史陶 2.186) represent a developmental stage between the single, bored round hollows and the double-linked ones (\pinyinterm{renwen} shakubun, introduction, pp. 117–122; cf. Itō [1962a], pp. 257-261; \pinyinterm{xu-jinxiong} [1974], pp. 20-21). For evidence that hollow shapes continued to evolve in Western \pinyinterm{zhou-dynasty} and Warring States times, see \pinyinterm{dong-short} (1956), p. 272. The triple-linked hollows found on scapulas in the upper stratum at 成-tzu-yai (\pinyinterm{li-ji} [1934], pp. 86, 88, fig. 9.4-6), if they are indeed of \pinyinterm{zhou-dynasty} date, represent a continuation of pre-\pinyinterm{anyang} practice (see n. 163).} 147
Pyromantic bones and shells excavated from the \pinyinterm{renmin-gongyuan} phase at \pinyinterm{zhengzhou}\footnote[148]{An 知-min (1954a), p. 34, fig. 3; KKHP (1958.3), p. 60, fig. 16.4, plate 7.6.}148
(which is thought to stand chronologically just prior to period I of the \pinyinterm{anyang} phase)\footnote[149]{长 光 (1968), p. 231.} 149
contain the first examples of double-hollows, a technical advance whose routine use parallels
the appearance of turtle shells as a divining medium.\footnote[150]{Itō (1962a), p. 260. On the nature and function of the double hollows, \ref{sec:chapters-ch01:1.5} (see see sec. 1.5). \pinyinterm{liu-yuanlin} (1974), pp. 113-114, provides an interesting account of six \pinyinterm{anyang} scapula fragments which, on the basis of their crude preparation, unusual hollow shapes, and writing style, he would date to a pre-period I stage of divination. The fragments (for photographs see his plate 12.1-6) are 3.2.139; 5.2.66; \pinyinterm{jiabian} 2342; 2344; 2875; \pinyinterm{yibian} 9105 (this last being from \pinyinterm{hougang}). In some cases, the hollows (all single) were burned on both shoulders, resulting in + -shaped cracks, but in other cases the heat source was placed so as to form a definite pattern of pu cracks, e.g., hollows on the left of the scapula were all burned on their right shoulder, and vice versa. Double hollows were presumably developed to facilitate the appearance of such crack patterns.} 150 Double hollows are found on
virtually all the period I turtle shells and may be regarded as standard for later, inscribed
shells. Double hollows also appear on period I scapulas, but generally not on scapulas of
periods II to V.\footnote[151]{\pinyinterm{liu-yuanlin} (1974), p. 118. 刘 (p. 117, plate 12.1) notes a rare period V bone fragment (\pinyinterm{jiabian} 3689) with a double hollow. Cf. \pinyinterm{xu-jinxiong} (1973a), pp. 17, 93–94.}151 The presence of double hollows on scapulas, therefore, suggests a period
I date.\footnote[152]{It is particularly important in this connection to emphasize a point stressed by \pinyinterm{xu-jinxiong} (see \ref{ch:1} (see ch. 1), n. 83): many of the round hollows on the shoulders of long hollows were not chiseled or drilled, but were merely the result of heat chipping; the distinction, which is crucial to his system of periodization, may not be visible in a rubbing.} 152
It has long been thought that more subtle distinctions in hollow shape might be used
to periodize the inscribed bones and shells with more discrimination, but it is only recently
% source: scan 142, printed 124
that a comprehensive study of the problem, primarily as it applies to the scapulimantic
hollows, has been undertaken by \pinyinterm{xu-jinxiong}.\footnote[153]{Due perhaps to the greater difficulty of cracking shell rather than bone, or to the desire to have cracks more rigorously patterned on plastrons than on scapulas, hollows made on shell did not vary as markedly with time as those made on bone; they have not yet been found useful for classification (cf. \pinyinterm{xu-jinxiong} [1974], pp. 12, 41).}153 The topic-which has led 徐 to
write several articles and two book-length monographs, the last superseding all the rest\footnote[154]{See \pinyinterm{xu-jinxiong} (1970a), (1972), (1972a), (1973), (1973a); his dissertation (1974); Menzies: Volume I, pp. 3-8; Volume II, pp. xv-xix, 7-9. For earlier discussion of hollow shapes, see \pinyinterm{chen-mengjia-spaced} (1956), pp. 10-11; \pinyinterm{renwen} shakubun, introduction, pp. 117-122; and the other works cited by \pinyinterm{xu-jinxiong-tight} (1974), p. 280, n. 16. It has been suggested (\pinyinterm{renwen} 848, shakubun, p. 311, n.; Itō [1962a], p. 261) that the various hollowmaking traditions represent different divining traditions brought to the \pinyinterm{shang} court by diviners from various regions. But, with the exception of the diviner 成, whom \pinyinterm{xu-jinxiong-tight} (1974), pp. 86-87, suggests may be linked to certain fatter, less neat, hollows, which he dates to late period I, no link has yet been established between particular diviners and particular hollow shapes. It would also be fruitful to look for links between engravers and hollow makers. See too n. 159 and n. 162.} 154
―is technical and the conclusions still tentative. Table 30, which summarizes some of his
preliminary results, indicates the type of conclusions to which his researches are leading
him; but his work, which involves the measurement and classification of hundreds of
hollow shapes, must be read in its entirety.\footnote[155]{To give but one example of his conclusions in periods III and IV, the trend in hollow length was from long to short. Hollows ``with wide, bowed shoulders and narrow rounded ends of giant size'' were gradually replaced by ``those with slightly curved shoulders and rounded ends of big size,'' which were replaced in turn by ``those with narrow shoulders and flat or triangular ends'' or by ``hollows of teardrop shape or middle or small size with one sharp and one rounded end. Each style comes at a particular time, but they overlap...'' (\pinyinterm{xu-jinxiong} [1974], p. 256).} 155 徐 is continuing to develop the analytical
vocabulary that this kind of study requires;\footnote[156]{Confronted with statements like ``the singed area is very large,'' one hollow ``is the same shape'' as another, hollows of one period ``are the same'' as those of another (\pinyinterm{xu-jinxiong} [1973], pp. 25, 123; 21, 108; 22, 111), one can only wonder, how large is large? Large in relation to what? To area of bone, area of unburned portion? And what constitutes “sameness''? These are not quibbles; I frequently find myself uncertain as to whether one hollow was, as 徐 claims, the same as another. \pinyinterm{xu-jinxiong-tight} (1974), p. iii, has now refined his terms: he classifies the hollows into eight types: straight shoulders and pointed ends; slightly curved shoulders and sharp ends; slightly curved shoulders and rounded ends with needlelike extension; slightly curved shoulders and rounded ends; bowed shoulders and narrow rounded ends; slightly curved shoulders and narrow rounded ends; straight shoulders and triangular ends; straight shoulders and tridentshaped ends. He is also paying greater attention to the techniques involved in making the hollows (e.g., pp. 58, 92).} 156 he is suggesting ways in which hollow shapes,
considered with other criteria such as writing style, hollow placement, and preface formula,
may permit us not only to separate period III inscriptions from those of period IV, but also
to classify some bones as belonging to early, middle, or late period III or early, middle,
or late period IV.\footnote[157]{\pinyinterm{xu-jinxiong} (1974), pp. 207-257; see too n. 155.}157 He is extending the size of his research sample;\footnote[158]{\pinyinterm{xu-jinxiong} (1973a) was based primarily on bones in the Royal Ontario Museum (Menzies, White) and the Carnegie Museum (\pinyinterm{kufang}) in Pittsburgh. His dissertation \listitem{1974} includes material held in Kyoto (\pinyinterm{renwen}) and at the Academia Sinica in Taiwan (\pinyinterm{jiabian} and \pinyinterm{yibian}). The revised version of the dissertation will also include materials from the British Museum (\pinyinterm{kufang}) and Cambridge University (京). 125 ]} 158 he is developing a
% source: scan 143, printed 125
technical rationale for some of the shape changes;\footnote[159]{See, for example, his attempts to relate hollow shape to the question of omen manipulation, or to the way in which certain shapes might facilitate cracking (\pinyinterm{xu-jinxiong} [1974], pp. 58, 83, 92, 97; cf. [1973], p. 126). Throughout he is aware that the careers of hollow makers did not necessarily correspond with those of diviners or engravers and is thus ready to assign transitional dates such as ``late I or early II'' when needed (e.g., [1974], pp. 86, 170, 207-208; cf [1973], pp. 6, 63).}\footnote[160]{\pinyinterm{xu-jinxiong} (1973a) and \listitem{1974} reproduce the hollows in detail; the glosses to White also describe hollow size and shape; the second volume of Menzies includes descriptions of the chiseled hollows.}\footnote[161]{My main hesitation about the reliability of \pinyinposs{xu-jinxiong} conclusions concerns the need to provide a sufficient number of touchstone inscriptions, dated by their ancestral titles (\ref{ch:4} (see ch. 4), n. 1), which would unequivocally tie certain hollow shapes to certain periods. His attempt (1973a, passim) to date the RFG inscriptions to period IV rather than period I, for example, is weakened by a circularity of argument in which elements of what he is attempting to prove are assumed to be true. He has not yet demonstrated, to my satisfaction at least, that hollows shapes which he finds characteristic of period IV should be dated to period IV in the first place (cf. my comments at n. 62). The matter will be discussed in detail in Studies. For a review of \pinyinterm{xu-jinxiong} (1973a) which makes related points, see Mickel, Journal of Oriental Studies 12 (1975), pp. 113-114.}\footnote[162]{It would be interesting, for example, to see whether bones and shells (such as those listed in appendix 3, n. 11) which may have been prepared before they were sent in to the \pinyinterm{shang} (\ref{ch:1} (see ch. 1), n. 69) reveal hollow-making traditions different from those that the \pinyinterm{shang} themselves practiced.}\footnote[163]{\pinyinterm{liu-yuanlin} (1974), pp. 99-112. It is true that three hollows were bored in a straight line on scapulimantic fragments excavated from the middle stratum at \pinyinterm{gaohuangmiao} so that they overlapped in what modern scholars have called the san-lien-tsuan ``triple-connected boring'' pattern (KKHP [1958.4], plate 8.1; \pinyinterm{liu-yuanlin} [1974], p. 109, plate 11). This arrangement, if intended to facilitate the cracking of the bone, was partly misguided; the linked hollows were bored horizontally and thus ran counter to the veins of the bone. The \pinyinterm{anyang} bones, by contrast, have hollows which are elongated vertically and thus run with the veins, producing cracks more readily (see \pinyinterm{zhang-bingquan} [1954], pp. 216, 217, fig. 9.4-6; \pinyinterm{li-ji} [1956], pp. 148-149; Itō [1962a], pp. 257, 258, fig. 26). This modification may have been prefigured by the san-lien-tsuan hollows on the back of 3.2.139, which may date from the pre-period I stage at \pinyinterm{anyang} (see n. 150).}\footnote[164]{See, for example, the scapulimantic remains from 刘里-ke and \pinyinterm{erligang}, conveniently reproduced by \pinyinterm{liu-yuanlin} (1974), plates 6 and 7.}
% source: scan 144, printed 127
scapulas from the \pinyinterm{renmin-gongyuan} phase at \pinyinterm{zhengzhou},\footnote[165]{See n. 148 above. I find no evidence to support the contention of Itō (1962a), p. 259, that on the bones and shells excavated from the \pinyinterm{renmin-gongyuan} phase, the bored circular hollow might be to either the right or left of the chiseled oval, resulting in no consistent direction to the cracks found on any one piece.} 165 and which was characteristic
of the hollows on the \pinyinterm{anyang} bones and shells, seems to have been related to the development of the double hollows and the formation of regular crack patterns.\footnote[166]{Cf. n. 150.}166

\subsubsection[Hollow Placement]{\origsecnum{4.3.2.4}Hollow Placement}
\label{sec:chapters-ch04:4.3.2.4}

The general placement of the hollows on bone and shell has already been described
for the historical period (sec. 1.5.1). \pinyinterm{xu-jinxiong} has discerned several patterns involving the long chiseled hollows on scapulas which, together with other criteria like the
hollow sizes and shapes, crack notations, and preface formulas, permit him to propose
refined dates for period III and period IV scapulas. The validity of his hypotheses remains
to be tested by further research.\footnote[167]{\pinyinterm{xu-jinxiong} (1974), p. 209, identifies three hollow patterns near scapula sockets: the ``one-one'' pattern. in which the two columns of hollows start at the same level (e.g., the back of \pinyinterm{jiabian} 446; 514: 531; \pinyinterm{renwen} 1863; 1994; 1999; Menzies 1670; 1673; 1674, etc.; these and the examples listed below are taken from the list at 徐 [1974], pp. 212-240); the ``one-two'' pattern, in which the inner column of hollows (i.e., that on the side of the socket notch) starts one space below the outer column (e.g., the back of \pinyinterm{jiabian} 546; 602; 638; \pinyinterm{renwen} 2054; 2117; 2299; Menzies 1856; 1887; 1999); and the ``one-three pattern, in which the inner column of hollows starts two spaces below the outer column (e.g., the back of \pinyinterm{jiabian} 669–671; \pinyinterm{renwen} 2189; 2263; 2355; Menzies 1672; 1677; 1697). He finds (p. 167) that the one-three pattern appears in periods III and IV; it may appear in V, but he has never found it in period I. He also finds (pp. 241-242) that the one-three pattern was normal in period III, whereas the one-one pattern was normal in period IV; and he argues for a situation in which the use of the one-three type gradually decreased and the use of the one-one type, which was related to the increasing use of shorter hollows, gradually increased (p. 244). The matter is complicated, however, by the fact that the one-one type was normal for the period III inscriptions which bear diviners' names (p. 245 thus he must posit a trend from one-one early in period III to onethree later in III, and back to one-one in IV (p. 245). Other explanations may be possible. In particular, greater rigor is needed to link specific hollow patterns to specific periods (cf. nn. 44 and 161). And it may also be necessary to relate hollow pattern (and size) to the size of the scapulas themselves, i.e., to the amount of working space available to the hollow maker.}167
\subsection*{Burn Marks}
The size and placement of the burn marks may also serve as periodizing criteria.
The application of heat in Neolithic and early \pinyinterm{shang} times was relatively primitive. Even
when hollows were prepared, the placing of the heat source and the resultant crack forms
might still be irregular.\footnote[168]{The nature of the Neolithic burn marks and cracks has not been well reported. \pinyinterm{liu-yuanlin} (1974), pp. 99-119, summarizes the evidence. Some of the \pinyinterm{erligang} scapulas were burned directly on the unprepared surface and also, on the same bone, in prepared hollows (\pinyinterm{zhengzhou} \pinyinterm{erligang}, p. 37). See n. 150 on what may have been early \pinyinterm{anyang} experiments with placing the heat source. In this connection, it would be worth studying the point at which the Neolithic or \pinyinterm{shang} diviners succeeded in producing pu cracks at all. The reports do not always provide specific information, but in some cases (e.g., the sheep scapulas from the 龙 site at Keshengzhuang E; see 豐西發掘報告, p. 68, plate 35) no cracks appear to have formed. 127 U ] B C}168
\pinyinterm{xu-jinxiong} is the first scholar to study the question intensively for the historical
period. He has noted that the absence of round chiseled hollows, which are not generally
% source: scan 145, printed 128
found on scapulas after period I (sec. 4.3.2.3), meant that more heat was required to crack
the bone; as a result, the area burned in and around the single, oval hollow was larger and
the bone surface was more frequently chipped in the later periods than it had been in
period\footnote[169]{\pinyinterm{xu-jinxiong} (1973a), pp. 19, 101.} 1.169 He has also discerned differences in the burn marks of periods III and IV.\footnote[170]{In period III, the crack was made by touching the bone lightly with the heat source to leave a large, perfectly formed burn mark; e.g., the back of Menzies 1775; 1812; 1933; 2068. In period IV, by contrast, a very hard, perhaps metal, object seems to have been pressed firmly into the bone; the burn mark was always small and tended to chip the bone more deeply than had been the case in period III; e.g., the back of Menzies 2343; 2462 for these and the previous examples, see the drawings in the front of 徐 [1973]); 2528 (photo no. 26 at the back of 徐 [1973]). The distinction ``is very easy to recognize from observation of the bones'' (徐 [1974], pp. 255-256). 徐 also notes ([1973], pp. 24, 119) that the heat source often touched both walls of a hollow in period IV, something that rarely happened in period III when the hollows were wider. For 徐's comments on the period III and IV hollow shapes, which must be considered jointly with the burn marks, see n. 155.} 170
The distinctions he draws are real, but his hypotheses about the way the burn marks may
be related to specific periods need further study.\footnote[171]{Cf. the questions raised in n. 161.} 171 In particular, such study will require
not only that hollow shapes on the back of bones be reproduced but that they be reproduced
in such a way that the character of the burned area can readily be identified.\footnote[172]{See, e.g., the symbols given by \pinyinterm{xu-jinxiong} (1973a), pp. 40-41; cf. n. 160.} 172
It will be clear from\ref{sec:chapters-ch04:4.3.2} (see  sec. 4.3.2) that \pinyinterm{xu-jinxiong}'s pioneering analyses of related
physical criteria--hollow shapes, hollow placement, burn marks—are helping to develop
significant new periodizing techniques. Future research will, it is hoped, refine the criteria
still further and apply them to an ever larger corpus of inscriptions, will critically evaluate
the way these criteria can be related to the inscriptions of particular periods, and will
provide a rationale for such stylistic changes. Physical criteria are of major importance
and must be considered by every scholar seeking to date the oracle bones.

\subsection*{Archaeological Criteria}

No detailed, comprehensive account of the relationship between oracle inscriptions and excavation sites has yet been published.\footnote[173]{For the individual archaeological reports, see \pinyinterm{guo-baojun} \listitem{1933} and (1933a); \pinyinterm{li-ji} (1930), (1933), (1949); \pinyinterm{shi-zhangru} (1933), (1945), (1947); \pinyinterm{dong-short} (1929a), (1929b), (1929c), (1933), (1936), (1949), (1949a). For summaries and interpretations, see Kaizuka and Itō (1953), pp. 19-30; \pinyinterm{hu-houxuan} (1955), pp. 43-121; \pinyinterm{chen-mengjia} (1956), pp. 35-42; Ito (1959); \pinyinterm{shi-zhangru} (1959), pp. 319-329; \pinyinterm{zhang-bingquan} (1967), pp. 831-851; \pinyinterm{dong-short} (1965), pp. 14-15; \pinyinterm{li-ji} (1977). It is valuable to read the pieces by Itō and \pinyinterm{shi-zhangru} in conjunction.} 173 And for the reasons given below, it is
not certain that one ever will be. Nevertheless, it is important to study, when possible, the
pit or trench from which particular oracle bones were excavated and their position with
regard to one another and other datable objects in the ground.\footnote[174]{This is only possible in the case of oracle bones that were scientifically excavated; for these, see the archaeological reports listed in n. 173. For an interesting attempt to deduce the location where certain privately excavated inscriptions must have been found, see Itō (1971), p. 88. 科} 174 Archaeological context
can serve as a primary criterion in providing a terminus post quem for the absolute date of
% source: scan 146, printed 129
an inscription; thus, an oracle bone found in an undisturbed Western \pinyinterm{zhou-dynasty} stratum is
likely to be no later than Western \pinyinterm{zhou-dynasty} in date; but the Western \pinyinterm{zhou-dynasty} context cannot
determine the terminus ante quem for the inscription, which might have been incised centuries
before.
Pit Provenance
For establishing relative dates within the \pinyinterm{shang} dynasty, pit provenance is a
secondary criterion whose usefulness may vary from pit to pit. This is so for a number of
reasons. First, if the contents of a pit are judged to be chronologically homogeneous, that
judgment-barring the presence of criteria such as ancestral titles or diviners' names on
every single fragment—can never be more than probable; the fact that one fragment of
unknown date was found with ninety-nine fragments dated to period I, say, does not
necessarily require that the unknown also be dated to period I, probable though that may
be.\footnote[175]{Pit E16 is a case in point. It contained some 300 fragments of pyromantic bone and shell; 6 pieces bore the names of RFG diviners (楚 : \pinyinterm{jiabian} 3oo3; Shao ): \pinyinterm{jiabian} 2941; 3012; 3045+3047; Tui Ẻ : \pinyinterm{jiabian} 3045+3047;3083; 甲图'u 163); all the other fragments have been dated to period I (e.g., \pinyinterm{jiabian} 2945; 2949; 2956; 2967) (董 [1933], p. 361, argued, on the basis of ambiguous ancestral titles, that \pinyinterm{jiabian} 2964, 2984, and 2989 should be dated to period II. Qu 万, in his k'ao-shih, dates these to period I. 董 also dated 4.0.0264 [containing the titles \pinyinterm{zujia} and \pinyinterm{zugeng}] to period II, but the registration number appears to be wrong; \pinyinterm{jiabian} 3165 [the present 4.0.0264] is a period I fragment containing neither of these titles.) 董 concluded that the pit contained divinations from periods I, II, and IVb (i.e., the RFG inscriptions). It would be more plausible, however, to consider that the RFG inscriptions should also be dated to period I (or II) (cf. \ref{ch:2} (see ch. 2), n. 18). But the contents of the pit, by themselves, cannot resolve the issue.} 175 Second, if the contents of a pit are chronologically mixed, it does not necessarily
follow that the oracle bones from the bottom stratum (assuming that the strata within
the pit have been accurately reported; see sec. 4.3.3.2) are earlier in date than those from
the top. If the oracle bones were first stored in sequence above ground, the earliest bones,
then on the bottom of the pile, would have been shifted to the top if the pile were inverted
when it was buried; and if the bones were poured into the pit,\footnote[176]{The slope of the shell fragments in YH127 indicated that they had been poured into the pit from the north (\pinyinterm{shi-zhangru} [1947], pp. 4142).} 176 or if there were more than
one pile above ground, the chronological relationships might have become thoroughly
confused.\footnote[177]{To cite but one example, all but one of the six RFG pieces (originally in eight fragments) in E16 (n. 175) were found at a depth of 8.0 to 8.8 meters, the bottom level of the pit. (The stratigraphy of E16 is given by \pinyinterm{li-ji} [1949], p. 251.) But the fragments of one RFG diviner, Shao, were found both at the bottom level (\pinyinterm{jiabian} 3012; 3045+3047) and at the top level (\pinyinterm{jiabian} 2941; at a depth of 3.5 to 4.5 meters). More marked chronological confusion was found in other pits; e.g., A26, E9 (董 [1929], p. 180); 5: H20; longitudinal trench 13 ping; TLK south sector (\pinyinterm{shi-zhangru} [1959], pp. 320, 321).} 177 All cases, in short, must be judged individually. Relative dates derived from
pit provenance are probabilities rather than certainties.\footnote[178]{The situation in particular pits, as well as the reasons why the \pinyinterm{shang} buried their oracle bones, will be considered in Studies. 129 1} 178
It is rarely possible to use other objects found in a pit, such as bronzes or pots, to date
175. Pit E16 is a case in point. It contained
some 300 fragments of pyromantic bone and shell;
pieces bore the names of RFG diviners (楚 :
\pinyinterm{jiabian} 3oo3; Shao ): \pinyinterm{jiabian} 2941; 3012;
3045+3047; Tui Ẻ : \pinyinterm{jiabian} 3045+3047;3083;
甲图'u 163); all the other fragments have been
dated to period I (e.g., \pinyinterm{jiabian} 2945; 2949;
2956; 2967) (董 [1933], p. 361, argued, on the
basis of ambiguous ancestral titles, that \pinyinterm{jiabian}
2964, 2984, and 2989 should be dated to period II.
Qu 万, in his k'ao-shih, dates these to period
I. 董 also dated 4.0.0264 [containing the titles
\pinyinterm{zujia} and \pinyinterm{zugeng}] to period II, but the
registration number appears to be wrong; \pinyinterm{jiabian} 3165 [the present 4.0.0264] is a period I
fragment containing neither of these titles.) 董
concluded that the pit contained divinations from
periods I, II, and IVb (i.e., the RFG inscriptions).
It would be more plausible, however, to consider
that the RFG inscriptions should also be dated to
period I (or II) (cf. \ref{ch:2} (see ch. 2), n. 18). But the contents
of the pit, by themselves, cannot resolve the issue.
176. The slope of the shell fragments in YH127
indicated that they had been poured into the pit
from the north (\pinyinterm{shi-zhangru} [1947], pp. 41-
42).
177. To cite but one example, all but one of
the six RFG pieces (originally in eight fragments)
in E16 (n. 175) were found at a depth of 8.0 to 8.8
meters, the bottom level of the pit. (The stratigraphy of E16 is given by \pinyinterm{li-ji} [1949], p. 251.)
But the fragments of one RFG diviner, Shao,
were found both at the bottom level (\pinyinterm{jiabian}
3012; 3045+3047) and at the top level (\pinyinterm{jiabian}
2941; at a depth of 3.5 to 4.5 meters). More
marked chronological confusion was found in other
pits; e.g., A26, E9 (董 [1929], p. 180);
5: H20; longitudinal trench 13 ping; TLK south
sector (\pinyinterm{shi-zhangru} [1959], pp. 320, 321).
178. The situation in particular pits, as well as
the reasons why the \pinyinterm{shang} buried their oracle
bones, will be considered in Studies.
% source: scan 147, printed 130
the inscriptions; the inscriptions, whatever their problems, provide the most precise clues
about date that we have, and they are usually used to date the other contents of the pit,
or the building foundations above it, rather than vice versa.¹ 179
Excavation Reports
Due in some cases to the inadequacy of the published reports and in other cases
to the absence of any report whatever, we are not as well informed as we should be about
the provenance of every ``scientifically excavated'' inscription. Excavators on the first
three \pinyinterm{anyang} digs (1928-1929), believing that floods had disordered the whole site and,
further, that the site had been occupied only during the short period from Kings \pinyinterm{wu-yi}
(period IVa) to Ti 乙 (period Va), did not pay full attention to the exact provenance of
the oracle bones, believing periodization both impossible and unnecessary. It was only at
the time of the fifth dig (fall 1931) that these preconceptions were swept away and the
excavators realized that the location of every bone and shell fragment was potentially
significant.\footnote[179]{Cf. Itō (1959); \pinyinterm{shi-zhangru} (1959), pp. 319-320. An exception to this generalization may be provided by the oracle bones excavated on the southern edge of \pinyinterm{xiaotun} in 1973. According to the preliminary report (KK [1975.1], pp. 27-46), many of the oracle bones were found in strata that could be classified as early, middle, or late and were found in association with pots and pot fragments, for which a stylistic and technical evolution could be developed. The fact that inscriptions of the Tui-diviner group from T53 were found with vessels of early rather than late style clearly has a bearing on the Old and New School controversy (see \ref{ch:2} (see ch. 2), n. 18). Hsiao Nan (1976), for example, uses the new evidence, both inscriptional and archaeological, to argue that the RFG inscriptions belong to late period I, but the validity of the pottery typology needs to be demonstrated more fully (as it will be, presumably, in the final report) before these conclusions can be accepted.}\footnote[180]{\pinyinterm{dong-short} (1933), p. 351.}\footnote[181]{The title of this projected work was 甲骨 wen-tzu yü 殷墟 yi-chih PTXFEKH ; it is referred to by \pinyinterm{dong-short} (1949), p. 240. Cf. n. 183.}\footnote[182]{A detailed report of digs 1 and 2 in sector F, on the south edge of the village, has not been published, possibly because no major ruins were found there (Itō [1959], p. 339. Nor are some of the early reports from other sectors sufficiently clear. As a result, we cannot be certain whether \pinyinterm{jiabian} 1-109 and 297-352 came from Eg, E16, E17, or E18; whether \pinyinterm{jiabian} 110-178, 368-375, and 391 came from A26 or A25\%; I find no provenance cited for \pinyinterm{jiabian} 181 and 282-296 (cf. \pinyinterm{chen-mengjia} [1956]. p. 169). We are uncertain if \pinyinterm{jiabian} 448-489 came from west of F36 (south of the village) or from longitudinal trenches I and 2 in sector B (陳, op. cit., p. 143); we can only suppose that \pinyinterm{jiabian} 490928 came from sector F within the village (陳, op. cit., pp. 143, 147) but that \pinyinterm{jiabian} 508 came from latitudinal trenches I and 2 in B. Similarly, we can only speculate that \pinyinterm{jiabian} 3224 and 3234 came from sectors B and E; that \pinyinterm{jiabian} 3636 and 3638 came from E; that \pinyinterm{jiabian} 3483 came from trenches F1-4 (陳, op. cit., pp. 143-144). No detailed report has been published about the provenance of the inscriptions excavated in digs 8 and 9 (Itō [1959], p. 350); thus, we have no details about the provenance of many of the inscriptions from \pinyinterm{jiabian} 3691 to 3938. Other examples of this sort can be given.}
% source: scan 148, printed 131
great majority, when we know the pit provenance of a scientifically excavated inscription,
we are frequently not informed about the stratigraphy within the pit, and thus the relative
position of the inscribed bones and shells is not clear.\footnote[183]{\pinyinterm{dong-short} (1949), p. 229, indicated that tables showing the relationship between the registration numbers and strata would eventually be published. One hopes that the information has been preserved in his notes (cf. n. 181 above) and that \pinyinterm{shi-zhangru}'s forthcoming \pinyinterm{xiaotun} ti-yi-pen: vichih ti fa-hsien yü fa-chüeh: ting pien : to be published by the Academia Sinica, will contain more information.} 183 And even when the stratigraphic
data are supplied, the system of registration numbers adopted by the archaeologists results
in inevitable ambiguities.\footnote[184]{For a description of this system see \ref{ch:1} (see ch. 1), n. 31. The stratigraphic ambiguity in the registration system arises from the fact that the second number arbitrarily divides originally contiguous materials into two numbered sequences of bone and shell. The third number may thus indicate the relative position of bone to bone or shell to shell, but it apparently fails to indicate the relative position of bone to shell. Thus, to give an example, we know that TLK, south sector, produced shells 3.0.0777-3.0.1000 and bones 3.2.0383-3.2.0923 (\pinyinterm{shi-zhangru} [1959], p. 321). If we assume (see below) that the sequence of the numbers reflects the situation in the pit, we can conclude that shell 3.0.0777 lay above 3.0.0877 which in turn lay above 3.0.0977 and that bone 3.2.0383 lay above 3.2.0483. But we have no way of telling from the registration numbers which shell lay above which bone and vice versa, for the registration system fails to indicate whether 3.0.0777, for example, lay above or below 3.2.0777. Yet another difficulty concerns the question of when the registration numbers were first assigned to the fragments; this is nowhere stated explicitly, but since, in more recent times, members of the \pinyinterm{anyang} Archaeological Research Station do not number pottery fragments until they reach the workroom (this information was obtained during the visit of the U.S. Paleoanthropology Delegation to \pinyinterm{anyang} on 1 June 1975), it is likely that the registration numbers were not assigned to the oracle-bone fragments till well after they had left the pit. It is not certain in each case, therefore, that the sequence of registration numbers finally assigned reflects the original lie of the bones in the pit (cf. n. 186). And in some instances, during the first dig \listitem{1928} at least, the registration numbers of oracle bones from different pits that were excavated simultaneously were commingled; e.g., 1.2.0123 (from A26; 20 October); 1.2.01281.2.0139 (from F24; 22 October); 1.2.0145 (from A26 again; 23 October); 1.2.0157 (from F31; 23 October); 1.2.0189 (from F24 again: 25 October). I glean this information from \pinyinterm{dong-short} (1929a), pp. 180, 181. This means that the registration numbers themselves do not necessarily provide reliable indications for estimating the total number of fragments in a pit; 1.2.0123 and I.2.0145, for example, were found in A26, but we cannot conclude that the fragments with the intervening 21 numbers also came from A26; some of them, at least, came from F24. Finally, it is not certain if the large numbers of blank oracle bones excavated were assigned registration numbers at all; to cite but one example, we are told that 721 of the 746 fragments excavated from F37 lacked writing (董 [1929b], p. 182), but the only registration numbers recorded run from 1.2.0171 (\pinyinterm{jiabian} 415) to 1.2.0218 (\pinyinterm{jiabian} 445), representing only 48 fragments. Those estimates of total numbers of fragments excavated (as in appendix 3, sec. I, therefore, that depend upon the registration numbers for their fragment counts may not be fully accurate, at least when using data from the first dig.} 184
In addition to these and other defects,\footnote[185]{For other problems, see, for example, Itō (1959), PP. 344, 354-355, 368, n. 11. The archaeological reports give the impression that the excavators were learning on the job. For example, \pinyinterm{shi-zhangru} (1947), p. 11, admits that the methods used in dig 14 were too hasty; and he notes that it was only in dig 15, i.e., after the great bulk of the inscriptions had been excavated, that two big advances in technique were introduced: a pulley device to measure the depth of strata and a camera frame for taking sequential photos. And to extend the list of problems still further, the reports are not always consistent. For example, according to \pinyinterm{shi-zhangru} (1959), p. 322, shells 15.0.0001-15.0.0360 were excavated from YH251; according to \pinyinterm{dong-short} (1949a), p. 288, some of these -shells 15.0.328-15.0.374-were taken from YH253, not YH251. Since the pits were not close, the contradiction is serious. Similarly, according to \pinyinterm{dong-short} (1949a), p. 288, shells 13.0.628131 ] U} 185 the archaeological reports present us with a
cold trail. Most were written thirty or forty years ago, many by scholars no longer living.
% source: scan 149, printed 132
The original sites have been covered over and the land returned to agriculture. It is
doubtful, therefore, that a fully accurate archaeology of the inscriptions excavated before
will ever be written. More rigorous archaeological techniques will presumably be@@
applied to new finds. The archaeology of the large number of oracle-bone inscriptions
excavated in spring 1973, for example, will, it is hoped, be reported in detail, so that this
and subsequent discoveries may eventually help resolve some of our dating problems. 186
Site Distribution
It is worth stressing, in conclusion, that the \pinyinterm{xiaotun} site has not been thoroughly
excavated. In the late 1920s and the 1930s, the Academia Sinica teams concentrated
primarily on the sector northeast of the village (that is, between the village and the Huan
River) and the sector on the southern edge of the village (see fig. 28). Other areas in the
immediate vicinity of \pinyinterm{xiaotun} have not been excavated, and even sections of the central
site where the building foundations were found were not thoroughly probed.\footnote[187]{Only the 乙 Z and 并 areas were excavated comprehensively; the 甲 section of foundations was not thoroughly trenched, and the foundations of 甲 8, 9, 14, and 15 were not completely excavated (\pinyinterm{shi-zhangru} [1959], pp. 317, 318, 330).} 187 This raises
the possibility that major inscription caches may still lie in the environs of, and even under
the buildings in, modern \pinyinterm{xiaotun}. Any discussion of the distribution of inscriptions by
came from YH127; according to 史
% out-of-scope-heading: 13.0.17714 came from YH127; according to 史

\pinyinterm{shi-zhangru} (1959), p. 322, shells 13.0.17715-
should also be included. At one point
% out-of-scope-heading: 13.0.17756 should also be included. At one point

\pinyinterm{chen-mengjia-spaced} (1956), p. 156, appears to agree
with \pinyinterm{shi-zhangru}; but at another point in the
same book (p. 170), he says that \pinyinterm{yibian} 8502-8531
(= 13.0.17718-13.0.17756) came from YH6, as
\pinyinterm{dong-short} ([1949a], p. 288) had reported. Again, the
pits are not close. Further contradictions of this
sort may be noted.
186. See n. 179 for the way recent finds may
help resolve the dispute over the periodization of
the RFG inscriptions. The limitations of the
reports of the 1930s were partly those of contemporary archaeology anywhere in the world, partly
those imposed by the unsettled condition of North
China, which must have made scientific excavation
extremely difficult. For an account of the banditry
in the \pinyinterm{anyang} region, for example, see Creel
(1935), p. 178; (1937), pp. 30-31 and plate II;
Karlbeck (1957), passim. Control of the site was
so uncertain that between the third and fourth
Academia Sinica digs (December 1929 and March
1930) the \pinyintext{henan} provincial authorities excavated
some 3,600 fragments, probably in the region
between pits A26 and A29 (the best published in
\pinyinterm{wenlu}), and the villagers excavated another 100
(62 of which are now at Columbia University,
38 at Cambridge University); see \pinyinterm{dong-short} (1945),
pt. 2, ch. 6, pp. 6a-b; Jung Keng (1947), pp. 24,
27; \pinyinterm{li-yan-spaced} (1966), p. 3; (1970), pp. 257, 319-
320. Similarly, the Sino-Japanese war (董
[1964], p. 52, editor's note) may account for our
ignorance about the stratigraphy of the bones and
shells in YH127, the largest cache of oracle bones
ever excavated. Given the way in which the great
wadded ball of mud and plastrons, weighing over
three tons, was packed into a coffinlike box,
subjected to heavy rains, bounced around on
trucks and trains, and damaged in 徐-chou-all
this before the pieces were separated (\pinyinterm{shi-zhangru} [1947], p. 8; \pinyinterm{li-ji} (1977), p. 116), it is
unlikely that the registration numbers accurately
reflect the lie of the plastrons in the ground. And
we do not even know which shells were at the
top and which shells were at the bottom of
the pit. If the RFG inscriptions are really from
period IVb (see \ref{ch:2} (see ch. 2), n. 18), it is particularly
regrettable that we do not know how they lay
with regard to all the period I inscriptions.
% source: scan 150, printed 133
period (appendix 3 must consider that the comparatively blank periods in our chronology
may be related to the blank spaces on the excavation maps.[\textasciicircum{}188 \textasciicircum{}188]: That our present III + IV inscriptions were excavated principally in sector F, south of \pinyinterm{xiaotun}, raises the possibility that the inscriptions buried under the modern village, if they exist, might also be from this period. This view is supported by the fact that most of the 4,829 inscribed bone and shell fragments excavated on the southern edge of \pinyinterm{xiaotun} in 1973 are said 188 to be dated to period IV (KK [1975.1], p. 39). The use of modern probing techniques, such as resistivity surveying (Goodyear [1971], pp. 229236, 239-243) may eventually permit caches of bone and other materials to be located under the village itself without inconveniencing the inhabitants. j 133 B 0 ] U

\section[Archaeological Criteria]{\origsecnum{4.3.3}Archaeological Criteria}
\label{sec:chapters-ch04:4.3.3}

No detailed, comprehensive account of the relationship between oracle inscriptions and excavation 
sites has yet been published.\footnote{For the individual archaeological reports, see \pinyinterm{guo-baojun} 
\listitem{1933} and (1933a); \pinyinterm{li-ji} (1930), (1933), (1949); \pinyinterm{shi-zhangru} (1933), (1945), (1947); \pinyinterm{dong-short} (1929a), 
(1929b), (1929c), (1933), (1936), (1949), (1949a). For summaries and interpretations, see Kaizuka and 
Itō (1953), pp. 19-30; \pinyinterm{hu-houxuan} (1955), pp. 43-121; \pinyinterm{chen-mengjia} (1956), pp. 35-42; Ito (1959); 
\pinyinterm{shi-zhangru} (1959), pp. 319-329; \pinyinterm{zhang-bingquan} (1967), pp. 831-851; \pinyinterm{dong-short} (1965), pp. 14-15; \pinyinterm{li-ji} (1977). It is valuable to read the pieces by Itō and \pinyinterm{shi-zhangru} in conjunction.} And for the reasons given below, 
it is not certain that one ever will be. Nevertheless, it is important to study, when possible, the pit or 
trench from which particular oracle bones were excavated and their position with regard to one another and 
other datable objects in the ground.\footnote{This is only possible in the case of oracle bones that were 
scientifically excavated; for these, see the archaeological reports listed in n. 173. For an interesting 
attempt to deduce the location where certain privately excavated inscriptions must have been found, see Itō 
(1971), p. 88.} Archaeological context can serve as a primary criterion in providing a terminus post quem 
for the absolute date of an inscription; thus, an oracle bone found in an undisturbed Western \pinyinterm{zhou-dynasty} stratum 
is likely to be no later than Western \pinyinterm{zhou-dynasty} in date; but the Western \pinyinterm{zhou-dynasty} context cannot determine the 
terminus ante quem for the inscription, which might have been incised centuries before.

\subsubsection[Pit Provenance]{\origsecnum{4.3.3.1}Pit Provenance}
\label{sec:chapters-ch04:4.3.3.1}

For establishing relative dates within the \pinyinterm{shang} dynasty, pit provenance is a secondary 
criterion whose usefulness may vary from pit to pit. This is so for a number of reasons. First, 
if the contents of a pit are judged to be chronologically homogeneous, that judgment---barring 
the presence of criteria such as ancestral titles or diviners' names on every single fragment---can 
never be more than probable; the fact that one fragment of unknown date was found with ninety-nine 
fragments dated to period I, say, does not necessarily require that the unknown also be dated to 
period I, probable though that may be.\footnote{Pit E16 is a case in point. It contained some 300 
fragments of pyromantic bone and shell; 6 pieces bore the names of RFG diviners (楚: \pinyinterm{jiabian} 3003; 
Shao: \pinyinterm{jiabian} 2941; 3012; 3045+3047; Tui: \pinyinterm{jiabian} 3045+3047;3083; 甲图'u 163); all the other 
fragments have been dated to period I (e.g., \pinyinterm{jiabian} 2945; 2949; 2956; 2967) (董 [1933], p. 361, 
argued, on the basis of ambiguous ancestral titles, that \pinyinterm{jiabian} 2964, 2984, and 2989 should be dated 
to period II. Qu 万, in his k'ao-shih, dates these to period I. 董 also dated 4.0.0264 [containing 
the titles \pinyinterm{zujia} and \pinyinterm{zugeng}] to period II, but the registration number appears to be wrong; \pinyinterm{jiabian} 
3165 [the present 4.0.0264] is a period I fragment containing neither of these titles.) 董 concluded that 
the pit contained divinations from periods I, II, and IVb (i.e., the RFG inscriptions). It would be more 
plausible, however, to consider that the RFG inscriptions should also be dated to period I (or II) (cf. \ref{ch:2} (see ch. 
2), n. 18). But the contents of the pit, by themselves, cannot resolve the issue.}

Second, if the contents of a pit are chronologically mixed, it does not necessarily follow that the oracle 
bones from the bottom stratum (assuming that the strata within the pit have been accurately reported; see 
sec. 4.3.3.2) are earlier in date than those from the top. If the oracle bones were first stored in sequence 
above ground, the earliest bones, then on the bottom of the pile, would have been shifted to the top if the 
pile were inverted when it was buried; and if the bones were poured into the pit,\footnote{The slope of the 
shell fragments in YH127 indicated that they had been poured into the pit from the north (\pinyinterm{shi-zhangru} [1947], 
pp. 41-42).} or if there were more than one pile above ground, the chronological relationships might have become 
thoroughly confused.\footnote{To cite but one example, all but one of the six RFG pieces (originally in eight 
fragments) in E16 (n. 175) were found at a depth of 8.0 to 8.8 meters, the bottom level of the pit. (The 
stratigraphy of E16 is given by \pinyinterm{li-ji} [1949], p. 251.) But the fragments of one RFG diviner, Shao, were found 
both at the bottom level (\pinyinterm{jiabian} 3012; 3045+3047) and at the top level (\pinyinterm{jiabian} 2941; at a depth of 3.5 
to 4.5 meters). More marked chronological confusion was found in other pits; e.g., A26, E9 (董 [1929], p. 180); 
5: H20; longitudinal trench 13 ping; TLK south sector (\pinyinterm{shi-zhangru} [1959], pp. 320, 321).}

All cases, in short, must be judged individually. Relative dates derived from pit provenance are probabilities 
rather than certainties.\footnote{The situation in particular pits, as well as the reasons why the \pinyinterm{shang} buried 
their oracle bones, will be considered in Studies.}
